\subsection{表示的特征标}

设 E 是 K 上有限维的左 A-模。E 的特征标或 E 的迹，记作 \(Tr_{E}\)，是线性型 \(a \mapsto \mathrm{Tr}(a_{\mathrm{E}})\)。设 \((e_{1}, \ldots, e_{n})\) 是 E 的一组基，\((e_{1}^{*}, \ldots, e_{n}^{*})\) 是它的对偶基。按定义，我们有关系式

\[
\mathrm{Tr} _ {\mathrm{E}} = \sum_ {i = 1} ^ {n} c _ {\mathrm{E}} (e _ {i}, e _ {i} ^ {*}).\tag{7}
\]

线性型 \(Tr_{E}\) 属于 \(\Theta_{\mathrm{E}}(\mathrm{A})\)。我们有 \(Tr_{E} = Tr_{E'}\)，这里 E 与 \(E'\) 是两个同构的 A-模。因此 \(Tr_{E}\) 只依赖于 E 的同构类。

设 \(\pi\) 是 A 在 K 上有限维向量空间 V 中的一个线性表示。\(\pi\) 的迹（有时也称 \(\pi\) 的特征标）定义为 A-模 \(V_{\pi}\) 的特征标，我们把它记作 \(\mathrm{Tr}_{\pi}\)。若 \((e_1, \ldots, e_n)\) 是 V 在 K 上的一组基，且 \((\pi_{ij}(a))\) 是 \(\pi(a)\) 关于这组基的矩阵，则我们有

\[
\operatorname{Tr} _ {\pi} (a) = \sum_ {i = 1} ^ {n} \pi_ {i i} (a).\tag{8}
\]

线性型 \(\mathrm{Tr}_{\pi}\) 属于 \(\Theta_{\pi}(\mathrm{A})\)。

命题 5. —— 设 S 是 K 上有限维的单 A-模，D 是 S 的交换子的反体，Z 是 D 的中心。下列性质等价：

\begin{enumerate}
\def\labelenumi{(\roman{enumi})}
\item
  S 的特征标 \(Tr_{S}\) 不为零。
\item
  存在元素 \(d \in \mathrm{D}\) 使得 \(\operatorname{Tr}_{\mathrm{D} / \mathrm{K}}(d) \neq 0\)。
\item
  \(\mathbf{Z}\) 是 \(\mathbf{K}\) 的可分扩张，且特征 \(p\)（\(\mathbf{K}\) 的特征）不整除次数 \([\mathrm{D}:\mathrm{Z}]\)。
\end{enumerate}

右 D-向量空间 S 是有限维的。设 \((e_{1},\ldots,e_{n})\) 是 S 在 D 上的一组基，u 是 \(\operatorname{End}_{\mathrm{D}}(\mathrm{S})\) 的一个元素。设 \((d_{ij})\) 是 u 关于基 \((e_{1},\ldots,e_{n})\) 的矩阵。我们有 \(u(e_{j})=\sum_{i=1}^{n}e_{i}d_{ij}\)（对 \(j\in[1,n]\)）。用 \(u_{K}\) 表示把 u 视为 K-向量空间 S 的自同态而得到的映射，用 \((u_{ij})\) 表示 \(u_{K}\) 关于 K-向量空间 S 的直和分解 \((e_{i}\mathrm{D})\) 的矩阵（II, §10, No.~5, 第 346 页）。我们有 \(\operatorname{Tr}(u_{\mathrm{K}})=\sum_{i}\operatorname{Tr}(u_{ii})\)。此外，\(u_{ii}\) 是 K-向量空间 \(e_{i}D\) 的自同态，它由 \(u_{ii}(e_{i}d)=e_{i}d_{ii}d\)（\(d\in D\)）定义，因此其迹等于 \(\operatorname{Tr}_{\mathrm{D}/\mathrm{K}}(d_{ii})\)。这样我们就证明了等式

\[
\mathrm{Tr} (u _ {\mathrm{K}}) = \mathrm{Tr} _ {\mathrm{D/K}} \bigl (\sum_ {i} d _ {i i} \bigr).\tag{9}
\]

由伯恩赛德定理（VIII, 第 83 页，命题 4 的推论 1），映射 \(a \mapsto a_{\mathrm{S}}\)（从 A 到 \(\mathrm{End}_{\mathrm{D}}(\mathrm{S})\) 的映射）是满射。性质 (i) 与 (ii) 的等价性于是由公式 (9) 得出。

若 K 的扩张 Z 是可分的，则存在元素 \(d \in Z\) 使得 \(\operatorname{Tr}_{\mathrm{Z/K}}(d) \neq 0\)（V, §8, No.~2, 第 49 页，命题 1 的推论）。此外，我们有 \(\operatorname{Tr}_{\mathrm{D/K}}(d) = [\mathrm{D} : \mathrm{Z}] \operatorname{Tr}_{\mathrm{Z/K}}(d)\)（III, §9, No.~4, 第 548 页，推论）；因此，若 p 不整除 {[}D : Z{]}，则我们有 \(\operatorname{Tr}_{\mathrm{D/K}}(d) \neq 0\)。这就证明了 (iii) 蕴含 (ii)。

若 \(\mathrm{Z}\) 不是 \(\mathrm{K}\) 的可分扩张，则 \(\operatorname{Tr}_{\mathrm{Z / K}}(x) = 0\)（对每个 \(x \in \mathrm{Z}\)），从而 \(\operatorname{Tr}_{\mathrm{D / K}}(d) = \operatorname{Tr}_{\mathrm{Z / K}}(\operatorname{Tr}_{\mathrm{D / Z}}(d)) = 0\)（对每个 \(d \in \mathrm{D}\)）。

现设 p 整除次数 \([D:Z]\)。于是它整除 D 在 Z 上的约化次数。对每个 \(d \in D\)，我们有 \(\mathrm{Tr}_{\mathrm{D}/\mathrm{Z}}(d) = 0\)（VIII, 第 344 页，注），从而 \(\mathrm{Tr}_{\mathrm{D}/\mathrm{K}}(d) = \mathrm{Tr}_{\mathrm{Z}/\mathrm{K}}(\mathrm{Tr}_{\mathrm{D}/\mathrm{Z}}(d)) = 0\)。这就完成了蕴含关系 (ii) \(\Rightarrow\) (iii) 的证明。

推论. —— 若体 \(\mathbf{K}\) 是完满的，则性质 (i)、(ii)、(iii) 均成立。

由于该体是完满的，K 的扩张 Z 是可分的（V, §15, No.~5, 第 125 页，定理 3）。性质 (iii) 于是由 VIII, 第 323 页的推论 3 得出。

命题 6. —— 设 \(\mathcal{S}_0\) 是 K 上有限维且迹非零的单 A-模的类所成的集合。线性型族 \((\mathrm{Tr}_{\mathrm{S}})_{\mathrm{S}\in \mathcal{S}_0}\) 在 K 上线性无关。

设 \(\mathrm{F}\) 是 \(\mathcal{S}_0\) 的一个有限子集，\(\mathrm{S}\) 是 \(\mathrm{F}\) 的一个元素。按假设，存在元素 \(a \in \mathrm{A}\) 使得 \(\operatorname{Tr}_{\mathrm{S}}(a) \neq 0\)。由命题 4 的推论 1（VIII, 第 83 页），存在元素 \(b \in \mathrm{A}\) 使得 \(b_{\mathrm{S}} = a_{\mathrm{S}}\)，且 \(b_{\mathrm{T}} = 0\)（对每个 \(\mathrm{T} \in \mathrm{F} - \{\mathrm{S}\}\)）。我们有 \(\operatorname{Tr}_{\mathrm{S}}(b) \neq 0\)，且 \(\operatorname{Tr}_{\mathrm{T}}(b) = 0\)（对 \(\mathrm{T} \in \mathrm{F} - \{\mathrm{S}\}\)）。于是族 \((\operatorname{Tr}_{\mathrm{S}})_{\mathrm{S} \in \mathrm{F}}\) 线性无关。命题 6 得证。

注. —— 设 \(S_{K}\) 是 K 上有限维单 A-模的类所成的集合。命题 6 也可由下述事实得到：诸 \(\Theta_{\mathrm{S}}(\mathrm{A})\)（S 取遍 \(S_{K}\)）之和是直和。

设

\[
0 \longrightarrow \mathrm{E} ^ {\prime} \longrightarrow \mathrm{E} \longrightarrow \mathrm{E} ^ {\prime \prime} \longrightarrow 0
\]

是 A-模的一个正合列，其中各模都在 K 上有限维。由 III, §9, No.~2, 第 543 页的命题 1，我们有 \(\mathrm{Tr}_{\mathrm{E}} = \mathrm{Tr}_{\mathrm{E}'} + \mathrm{Tr}_{\mathrm{E}''}\)。由格罗滕迪克群 \(\mathrm{R}_{\mathrm{K}}(\mathrm{A})\) 的定义（VIII, 第 194 页）及其泛性质（VIII, 第 186 页，命题 4），存在一个加群同态 \(\theta\)，从 \(\mathrm{R}_{\mathrm{K}}(\mathrm{A})\) 到 \(\mathrm{A}^*\)，它由关系式 \(\theta([\mathrm{E}]) = \mathrm{Tr}_{\mathrm{E}}\)（对每个 K 上有限维的 A-模 E 成立）刻画。特别地，E 的半单化的迹等于 E 的迹。由 \(\theta\) 我们导出一个 K-线性映射 \(\theta_{\mathrm{K}}: \mathrm{K} \otimes_{\mathbf{Z}} \mathrm{R}_{\mathrm{K}}(\mathrm{A}) \to \mathrm{A}^*\)。

推论. —— a) 设体 K 的特征为 0。同态 \(\theta\) 与 \(\theta_{K}\) 都是单射。两个 K 上有限维的半单 A-模同构，当且仅当它们的特征标相等。

\begin{enumerate}
\def\labelenumi{\alph{enumi})}
\setcounter{enumi}{1}
\tightlist
\item
  设体 K 是特征为非零素数 p 的完满体。那么同态 \(\theta_{K}\) 是单射，且 \(\theta\) 的核为 \(p\mathrm{R}_{\mathrm{K}}(\mathrm{A})\)。设 E 是 K 上有限维的半单 A-模。线性型 \(Tr_{E}\) 为零，当且仅当存在 A-模 F 使得 E 同构于 \(F^{p}\)。
\end{enumerate}

设 \(S_{K}\) 是 K 上有限维单 A-模的类所成的集合。诸元素 {[}S{]}（其中 S 取遍 \(S_{K}\)）构成 Z-模 \(\mathrm{R}_{\mathrm{K}}(\mathrm{A})\) 的一组基，因而诸元素 \(1 \otimes [S]\) 构成 K-向量空间 \(\mathrm{K} \otimes_{\mathbf{Z}} \mathrm{R}_{\mathrm{K}}(\mathrm{A})\) 的一组基（VIII, 第 195 页）。

设体 K 是完满的。由命题 6 以及 VIII, 第 383 页的推论，元素 \(\theta([S]) = \theta_{\mathrm{K}}(1 \otimes [S]) = \mathrm{Tr}_{\mathrm{S}}\)（\(A^{*}\) 中的元素）在 K 上线性无关。由此可知同态 \(\theta_{K}\) 是单射，且同态 \(\theta\) 的核由元素 \(\sum_{S \in S_{K}} n_{S}[S]\)（\(\mathrm{R}_{\mathrm{K}}(\mathrm{A})\) 的元素，满足 \(n_{S} \cdot 1_{K} = 0\) 对每个 \(S \in S_{K}\)）组成。因此 \(\theta\) 的核等于 \(p R_{\mathrm{K}}(\mathrm{A})\)，其中 p 是 K 的特征。特别地，若 K 的特征为 0，则同态 \(\theta\) 是单射。

a) 的最后一个断言由 VIII, 第 190 页的推论得出。

设 b) 的假设成立，并设 E 是 K 上有限维的半单 A-模。若存在 A-模 F 使 E 是 p 个同构于 F 的子模的直和，则模 F 是半单的且在 K 上有限维，并且我们有 \(Tr_{E} = p Tr_{F} = 0\)。反之，设 \(Tr_{E} = 0\)。我们有 \([E] \in p R_{K}(A)\)，故每个单模 \(S \in S_{K}\) 在 E 中的重数都是 p 的倍数（VIII, 第 190 页，命题 7）；我们把它写作 \(pn_{S}\)，其中 \(n_{S} \in N\)。族 ( \(n_{S}\) ) 具有有限支集。令 \(F = \oplus_{S \in S_{K}} S^{n_{S}}\)。于是我们有 {[}E{]} = {[}F\^{}\{p\}{]}，从而 E 与 F\^{}\{p\} 同构（VIII, 第 190 页，推论）。

设 \(\mathrm{E}\) 是 \(\mathrm{K}\) 上有限维的 A-模。设 \(a \in \mathrm{A}\)。用 \(\chi_{\mathrm{E}}(a; \mathrm{T})\) 表示 \(1_{\mathrm{E}} + T a_{\mathrm{E}}\)（\(\mathrm{K}[\mathrm{T}]\)-模 \(\mathrm{E}[\mathrm{T}] = \mathrm{K}[\mathrm{T}] \otimes_{\mathrm{K}} \mathrm{E}\) 的自同态）的行列式（III, §8, No.~11, 第 541 页）。我们有关系式

\[
\chi_ {\mathrm{E}} (a; \mathrm{T}) \equiv 1 + \mathrm{Tr} _ {\mathrm{E}} (a) \mathrm{T} \pmod {\mathrm{T} ^ {2} \mathrm{K} [ \mathrm{T} ]}\tag{10}
\]

（同前引文，公式 (49)）。由于这个多项式的常数项等于 1，它是一个可逆的形式幂级数（IV, §4, No.~4, 第 30 页）。此外，若

\[
0 \longrightarrow \mathrm{E} ^ {\prime} \longrightarrow \mathrm{E} \longrightarrow \mathrm{E} ^ {\prime \prime} \longrightarrow 0
\]

是 A-模的一个正合列，且各模都在 K 上有限维，则我们有 \(\chi_{\mathrm{E}}(a;\mathrm{T})=\chi_{\mathrm{E}^{\prime}}(a;\mathrm{T})\chi_{\mathrm{E}^{\prime\prime}}(a,\mathrm{T})\)（III, §8, No.~7, 第 534 页，公式 (31)）。由格罗滕迪克群 \(\mathrm{R}_{\mathrm{K}}(\mathrm{A})\) 的定义（VIII, 第 194 页）及其泛性质（VIII, 第 186 页，命题 4），存在唯一的同态 \(\chi_{a}\)，从群 \(\mathrm{R}_{\mathrm{K}}(\mathrm{A})\) 到乘法群 \(1+\mathrm{TK}[[\mathrm{T}]]\)，使得对每个 K 上有限维的 A-模 E 有 \(\chi_{a}([\mathrm{E}])=\chi_{\mathrm{E}}(a;\mathrm{T})\)。由公式 (10) 我们导出关系式

\[
\chi_ {a} (x) \equiv 1 + \theta (x) (a) \mathrm{T} \pmod {\mathrm{T} ^ {2} \mathrm{K} [ [ \mathrm{T} ] ]}\tag{11}
\]

其中 \(x\in \mathrm{R}_{\mathrm{K}}(\mathrm{A})\)，\(a\in \mathrm{A}\)。

若 E 与 \(E'\) 是两个 K 上有限维且半单化同构的 A-模，则我们有 \(\chi_{\mathrm{E}}(a;\mathrm{T})=\chi_{\mathrm{E}'}(a;\mathrm{T})\)。

定理 2. —— 设 \(\mathcal{A}\) 是 K-向量空间 A 的一个生成子集。同态 \(\chi_{\mathcal{A}}: \mathrm{R}_{\mathrm{K}}(\mathrm{A}) \to (1 + \mathrm{TK}[[\mathrm{T}]])^{\mathcal{A}}\) 由关系式 \(\chi_{\mathcal{A}}(x) = (\chi_a(x))_{a \in \mathcal{A}}\) 定义，它是单射。

设 \(x\) 是 \(\mathrm{R}_{\mathrm{K}}(\mathrm{A})\) 中满足 \(\chi_{\mathcal{A}}(x) = 1\) 的一个元素。由 (11)，我们有 \(\theta(x)(a) = 0\)（对每个 \(a \in \mathcal{A}\)），从而 \(\theta(x) = 0\)，因为 \(\theta(x)\) 是 A 上的 K-线性型而 \(\mathcal{A}\) 生成 K-向量空间 A。若 K 的特征为零，这就蕴含 \(x = 0\)（VIII, 第 384 页，命题 6 的推论），从而在这种情形下结论成立。以下设 K 的特征 \(p\) 非零。

我们先处理 K 是代数闭的情形。由前述及同前引文，此时我们有 \(x \in p \mathrm{R}_{\mathrm{K}}(\mathrm{A})\)。设 \(y \in \mathrm{R}_{\mathrm{K}}(\mathrm{A})\) 满足 x = py。对 A 的每个元素 a，我们有 \(\chi_{a}(y)^{p} = \chi_{a}(py) = \chi_{a}(x) = 1\)。于是 \((\chi_{a}(y) - 1)^{p} = 0\)，故 \(\chi_{a}(y) = 1\)，因为环 K{[}{[}T{]}{]} 是整环。所以 y 属于自同态 \(\chi_{A}\) 的核。由归纳法可知，x 属于 \(p^{n}\mathrm{R}_{\mathrm{K}}(\mathrm{A})\)（对每个整数 \(n \geqslant 1\)）。由于 \(\mathrm{R}_{\mathrm{K}}(\mathrm{A})\) 是自由 Z-模，这蕴含 x = 0，从而在这种情形下 \(\chi_{A}\) 是单射。

若不再假设 K 是代数闭的，我们取定 K 的一个代数闭包 \(\overline{K}\)，并考虑如下的群与群同态构成的图

\[
\begin{array}{c} \mathrm{R} _ {\mathrm{K}} (\mathrm{A}) \xrightarrow {\chi_ {\mathcal {A}}} (1 + \mathrm{TK} [ [ \mathrm{T} ] ]) ^ {\mathcal {A}} \\ \Biggl \downarrow^ {u} \qquad \qquad \qquad \qquad \Biggl \downarrow^ {i} \\ \mathrm{R} _ {\overline {{\mathrm{K}}}} (\mathrm{A} _ {(\overline {{\mathrm{K}}})}) \xrightarrow {\overline {{\chi}} _ {\mathcal {A}}} (1 + \mathrm{T} \overline {{\mathrm{K}}} [ [ \mathrm{T} ] ]) ^ {\mathcal {A}}, \end{array}\tag{12}
\]

其中 \(u\) 是由 K 到 \(\overline{\mathbf{K}}\) 的标量扩张导出的同态（VIII, 第 195 页），\(i\) 是典范单射，而 \(\overline{\chi}_{\mathcal{A}}\) 是同态 \(z \mapsto (\chi_{1 \otimes a}(z))_{a \in \mathcal{A}}\)。由 III, §9, No.~1, 第 542 页的公式 (12)，图 (12) 交换。由前述，同态 \(\overline{\chi}_{\mathcal{A}}\) 是单射。由于 \(u\) 是单射（VIII, 第 195 页，定理 1），同态 \(\chi_{\mathcal{A}}\) 是单射。

推论 1. —— 设 E 与 F 是 K 上有限维的半单 A-模，\(\mathcal{A}\) 是 K-向量空间 A 的一个生成子集。假设对每个 \(a \in \mathcal{A}\)，K-向量空间 E 与 F 的自同态 \(a_{E}\) 与 \(a_{F}\) 的特征多项式都相等。那么 A-模 E 与 F 同构。

设 \(a\) 是 \(\mathcal{A}\) 的一个元素。\(a_{\mathrm{E}}\) 与 \(a_{\mathrm{F}}\) 的特征多项式有相同的次数，故 E 的维数等于 F 的维数；我们把它记作 \(n\)。设 \(\mathrm{Pc}_{\mathrm{E}}(a;\mathrm{T})\) 是 \(a_{\mathrm{E}}\) 的特征多项式。在 \(\mathrm{K(T)}\) 中，我们有等式

\[
\chi_ {\mathrm{E}} (a; \mathrm{T}) = \det \left(1 + a _ {\mathrm{E}} \mathrm{T}\right) = (- \mathrm{T}) ^ {n} \det \left(\frac {- 1}{\mathrm{T}} - a _ {\mathrm{E}}\right) = (- \mathrm{T}) ^ {n} \operatorname{Pc} _ {\mathrm{E}} \left(a; \frac {- 1}{\mathrm{T}}\right),
\]

而 \(\chi_{\mathrm{F}}(a;\mathrm{T})\) 由一个类似的公式给出。由我们的假设，有 \(\chi_{\mathrm{E}}(a;\mathrm{T})=\chi_{\mathrm{F}}(a;\mathrm{T})\)。由定理 2，我们有 \([E]=[F]\)，这蕴含 E 与 F 同构（VIII, 第 190 页，命题 7 的推论）。

推论 2. —— 设 A 是 K 上有限次数的中心单代数。设 B 是半单 K-代数，f 与 g 是从 B 到 A 的代数同态，\(\mathcal{B}\) 是 K-向量空间 B 的一个生成子集。下列性质等价：

\begin{enumerate}
\def\labelenumi{(\roman{enumi})}
\item
  存在 A 的内自同构 \(\theta\) 使得 \(g = \theta \circ f\)。
\item
  对每个 \(b \in \mathcal{B}\)，有 \(\mathrm{Pc}_{\mathrm{A / K}}(f(b); \mathrm{X}) = \mathrm{Pc}_{\mathrm{A / K}}(g(b); \mathrm{X})\)。
\end{enumerate}

当 \(\mathrm{K}\) 的特征为零时，这些性质等价于下面的性质：

\begin{enumerate}
\def\labelenumi{(\roman{enumi})}
\setcounter{enumi}{2}
\tightlist
\item
  对每个 \(b \in \mathcal{B}\)，有 \(\mathrm{Tr}_{\mathrm{A / K}}(f(b)) = \mathrm{Tr}_{\mathrm{A / K}}(g(b))\)。
\end{enumerate}

把 A 的加群分别赋予外部合成律 \((b,a)\mapsto f(b)a\) 与 \((b,a)\mapsto g(b)a\)，将这样得到的左 B-模记作 M 与 N。由 VIII, 第 257 页的命题 1，性质 (i) 等价于 B-模 M 与 N 同构这一事实。按构造，与 B 的元素 b 相应的位似 \(b_{M}\) 是 A 中由 \(f(b)\) 作的左乘；因此我们有关系式

\[
\begin{array}{c} \mathrm {Pc_ {M}} (b; \mathrm{X}) = \mathrm {Pc_ {A / K}} (f (b); \mathrm{X}), \\ \mathrm{Tr} _ {\mathrm{M}} (b) = \mathrm{Tr} _ {\mathrm{A/K}} (f (b)) \end{array}
\]

以及把 M 换成 N、f 换成 g 后的两个类似关系式。我们知道（VIII, 第 386 页，推论 1），B-模 M 与 N 同构当且仅当 \(\mathrm{Pc}_{\mathrm{M}}(b;\mathrm{X})=\mathrm{Pc}_{\mathrm{N}}(b;\mathrm{X})\) 对每个 \(b\in B\) 成立。当体 K 的特征为 0 时，该关系式也等价于 \(\mathrm{Tr}_{\mathrm{M}}(b)=\mathrm{Tr}_{\mathrm{N}}(b)\) 对每个 \(b\in B\) 成立（VIII, 第 384 页，命题 6 的推论）。性质 (i) 与 (ii) 的等价性，以及当 K 的特征为 0 时 (i) 与 (iii) 的等价性，由此得出。

\subsection{模类集合的系数}

设 C 是 A-模的类的一个遗传集（VIII, 第 183 页，定义 1）。我们假设每个 C 型 A-模都在 K 上有限维，并用 \(\Theta_{\mathcal{C}}(\mathrm{A})\) 表示 C 型 A-模的系数所成的集合。由 VIII, 第 378 页的命题 2，集合 \(\Theta_{\mathcal{C}}(\mathrm{A})\) 也是诸 A-线性映射 \(u: E \to A^{*}\)（其中 E 取遍 C）的像的并；它也是子空间 \(\Theta_{\mathrm{E}}(\mathrm{A})\)（属于 \(A^{*}\)，其中 E 取遍 C）的并（同前引文）。(A, A)-子双模族 \((\Theta_{\mathrm{E}}(\mathrm{A}))_{\mathrm{E} \in \mathcal{C}}\)（\(A^{*}\) 的 (A, A)-子双模族）是有向的，故其并 \(\Theta_{\mathcal{C}}(\mathrm{A})\) 是 \(A^{*}\) 的一个 (A, A)-子双模。

命题 7. —— \(A^{*}\) 的一个 K 上有限维的左 A-子模包含于 \(\Theta_{\mathcal{C}}(\mathrm{A})\)，当且仅当它是 C 型的。

设 V 是 \(A^{*}\) 的一个 K 上有限维的左 A-子模。我们在 VIII, 第 379 页注 3 中已看到有 \(V \subset \Theta_{\mathrm{V}}(\mathrm{A})\)。若 V 是 C 型的，则 \(\Theta_{\mathrm{V}}(\mathrm{A}) \subset \Theta_{\mathcal{C}}(\mathrm{A})\)，故 V 包含于 \(\Theta_{\mathcal{C}}(\mathrm{A})\)。反之，设 V 包含于 \(\Theta_{\mathcal{C}}(\mathrm{A})\)。由于 \(\Theta_{\mathcal{C}}(\mathrm{A})\) 是有向族 \((\Theta_{\mathrm{E}}(\mathrm{A}))_{\mathrm{E}\in\mathcal{C}}\) 的并，且 V 在 K 上有限维，故存在一个 C 型模 E 使 V 包含于 \(\Theta_{\mathrm{E}}(\mathrm{A})\)。而存在自然数 n 使 \(\Theta_{\mathrm{E}}(\mathrm{A})\) 同构于 \(\mathrm{E}^{n}\) 的一个商模（VIII, 第 378 页，命题 2）。由于 C 是遗传的，V 是 C 型的。

推论. —— 设 E 是 K 上有限维的 A-模。那么 E 是 C 型的，当且仅当 \(\Theta_{\mathrm{E}}(\mathrm{A})\) 包含于 \(\Theta_{\mathcal{C}}(\mathrm{A})\)。

该条件显然是必要的。

反之，设 \(\Theta_{\mathcal{C}}(\mathrm{A})\) 包含 \(\Theta_{\mathrm{E}}(\mathrm{A})\)。A-模 \(\Theta_{\mathrm{E}}(\mathrm{A})\) 在 K 上有限维（VIII, 第 378 页，命题 2），故由命题 7 它是 C 型的。而存在整数 \(n \geqslant 0\) 使 E 同构于 \(\Theta_{\mathrm{E}}(\mathrm{A})^{n}\) 的一个 A-子模（VIII, 第 378 页，命题 2）。由于 C 是遗传的，E 是 C 型的。

\subsection{限制对偶上的余代数结构}

对任意的 \(a\)（\(\mathbf{A}\) 中的），我们用 \(\eta(a)\) 表示线性型 \(f \mapsto f(a)\)（\(\Theta(\mathbf{A})\) 上的）。这样我们就定义了一个 K-线性映射 \(\eta\)，从 \(\mathbf{A}\) 到对偶 \(\Theta(\mathbf{A})^*\)（向量空间 \(\Theta(\mathbf{A})\) 的对偶）。令 \(\varepsilon = \eta(1)\)。

命题 8. —— 在向量空间 \(\Theta(A)\) 上存在唯一的余代数结构（III, §11, No.~1, 第 574 页），使得映射 \(\eta: A \to \Theta(A)^*\) 是从 \(A\) 到 \(\Theta(A)\) 的对偶代数（III, §11, No.~2, 第 579 页）的同态。余代数 \(\Theta(A)\) 是余结合的，且以 \(\varepsilon\) 为余单位。

对每个整数 \(n \geqslant 1\)，考虑 K-线性映射 \(j_{n}\)（从 \(\Theta(\mathrm{A})^{\otimes n}\) 到 \(A^{\otimes n}\) 的对偶），它由公式

\[
\langle j _ {n} \left(f _ {1} \otimes \dots \otimes f _ {n}\right), a _ {1} \otimes \dots \otimes a _ {n} \rangle = \prod_ {i = 1} ^ {n} \left\langle f _ {i}, a _ {i} \right\rangle\tag{13}
\]

刻画，其中 \((a_{i})\) 属于 \(\mathrm{A}^n\)，\((f_i)\) 属于 \(\Theta (\mathrm{A})^n\)。由 II, §7, No.~7, 第 308 页的命题 16 (ii)，映射 \(j_{n}\) 是单射。用 \(m_{\mathrm{K}}:\mathrm{K}\otimes \mathrm{K}\to \mathrm{K}\) 与 \(m_{\mathrm{A}}:\mathrm{A}\otimes \mathrm{A}\rightarrow \mathrm{A}\) 分别表示由 K 与 A 中的乘法导出的映射。对 \(f,g\in \Theta (\mathrm{A})\) 与 \(a,b\in \mathrm{A}\)，我们有

\[
\langle j _ {2} (f \otimes g), a \otimes b \rangle = m _ {\mathrm{K}} \circ (\eta (a) \otimes \eta (b)) (f \otimes g).
\]

因此我们有

\[
\langle j _ {2} (t), a \otimes b \rangle = m _ {\mathrm{K}} \circ (\eta (a) \otimes \eta (b)) (t)\tag{14}
\]

对每个 \(t\in \Theta (\mathrm{A})\otimes \Theta (\mathrm{A})\) 成立。

引理 2. —— 设 \(c: \Theta(A) \to \Theta(A) \otimes \Theta(A)\) 是一个 K-线性映射。那么 \(\eta\) 是从 A 到余代数 \((\Theta(A), c)\) 的对偶代数的同态，当且仅当下面的图交换：

\[
\begin{array}{c} \Theta (\mathrm{A}) \xrightarrow {c} \Theta (\mathrm{A}) \otimes \Theta (\mathrm{A}) \\ \Biggl \downarrow_ {j _ {1}} \qquad \qquad \qquad \qquad \Biggl \downarrow_ {j _ {2}} \\ \mathrm{A} ^ {*} \xrightarrow {^ t m _ {\mathrm{A}}} (\mathrm{A} \otimes \mathrm{A}) ^ {*}. \end{array}\tag{15}
\]

事实上，\(\eta\) 是从 A 到余代数 \((\Theta (\mathrm{A}),c)\) 的对偶代数的同态，当且仅当 \(\eta (ab) = m_{\mathrm{K}}\circ (\eta (a)\otimes \eta (b))\circ c\) 对一切 \(a,b\in \mathrm{A}\) 成立，也就是

\[
\eta (a b) (f) = m _ {\mathrm{K}} \circ (\eta (a) \otimes \eta (b)) (c (f))
\]

对 \(a, b \in \mathrm{A}\) 与 \(f \in \Theta(\mathrm{A})\) 成立。而今，我们有

\[
\eta (a b) (f) = f (a b) = \left\langle {} ^ {t} m _ {\mathrm{A}} \left(j _ {1} (f)\right), a \otimes b \right\rangle
\]

而由 (14)，

\[
m _ {K} \circ (\eta (a) \otimes \eta (b)) (c (f)) = \langle j _ {2} (c (f)), a \otimes b \rangle
\]

对一切 \(a, b \in \mathrm{A}\) 与每个 \(f \in \Theta(\mathrm{A})\) 成立。引理得证。

由于 \(j_{2}\) 是单射，至多存在一个使上图交换的线性映射 c。为证明其存在性，必须证明 \(^{t}m \circ j_{1}\) 的像包含于 \(j_{2}\) 的像。换言之，必须证明：对 \(\Theta(\mathrm{A})\) 的每个元素 f，存在一个自然数 n 以及元素 \(f_{1}^{\prime}, \ldots, f_{n}^{\prime}, f_{1}^{\prime\prime}, \ldots, f_{n}^{\prime\prime}\)（\(\Theta(\mathrm{A})\) 中的元素），满足关系式

\[
f (a b) = \sum_ {i = 1} ^ {n} f _ {i} ^ {\prime} (a) f _ {i} ^ {\prime \prime} (b)\tag{16}
\]

其中 \(a, b \in \mathbf{A}\)。此时我们将有

\[
c (f) = \sum_ {i = 1} ^ {n} f _ {i} ^ {\prime} \otimes f _ {i} ^ {\prime \prime}.\tag{17}
\]

由 VIII, 第 379 页的推论，存在一个以 f 为系数的 K 上有限维左 A-模 E。设 \((e_{1},\ldots,e_{n})\) 是 E 的一组基，\((e_{1}^{*},\ldots,e_{n}^{*})\) 是对偶基，x 是 E 的一个元素，\(x^{*}\) 是 \(E^{*}\) 的一个元素且满足 \(f=c_{\mathrm{E}}(x,x^{*})\)。令 \(f_{i}^{\prime}=c_{\mathrm{E}}(e_{i},x^{*})\) 与 \(f_{i}^{\prime\prime}=c_{\mathrm{E}}(x,e_{i}^{*})\)（对 \(i\in[1,n]\)）；对

A 中的 \(a, b\)，我们有

\[
\begin{array}{r l} & {f (a b) = \langle x ^ {*}, a b x \rangle = \langle x ^ {*} a, b x \rangle = \left\langle \sum_ {i} \langle x ^ {*} a, e _ {i} \rangle e _ {i} ^ {*}, b x \right\rangle} \\ & {\qquad = \sum_ {i} \langle x ^ {*} a, e _ {i} \rangle \langle e _ {i} ^ {*}, b x \rangle = \sum_ {i} \langle x ^ {*}, a e _ {i} \rangle \langle e _ {i} ^ {*}, b x \rangle = \sum_ {i} f _ {i} ^ {\prime} (a) f _ {i} ^ {\prime \prime} (b),} \end{array}
\]

从而得到 (16)。

我们来证明 \(c\) 的余结合性。为此，考虑 K-线性映射

\[
c ^ {\prime} = \left(c \otimes 1 _ {\Theta (\mathrm{A})}\right) \circ c \quad \text { and } \quad c ^ {\prime \prime} = \left(1 _ {\Theta (\mathrm{A})} \otimes c\right) \circ c
\]

它们都是从 \(\Theta (\mathrm{A})\) 到 \(\Theta (\mathrm{A})^{\otimes 3}\) 的映射。我们有关系式

\[
\begin{array}{r l} \langle j _ {3} (f \otimes c (g)), a \otimes b \otimes c \rangle & = \langle f, a \rangle \langle j _ {2} \circ c (g), b \otimes c \rangle \\ & = \langle f, a \rangle \langle g, b c \rangle \\ & = \langle j _ {2} (f \otimes g), a \otimes b c \rangle \\ & = \langle^ {t} (\mathrm{Id} _ {\mathrm{A}} \otimes m _ {\mathrm{A}}) \circ j _ {2} (f \otimes g), a \otimes b \otimes c \rangle \end{array}
\]

其中 \(f, g \in \Theta(\mathbf{A})\)，\(a, b, c \in \mathbf{A}\)。由此我们导出下面的图是交换的：

\[
\begin{array}{c} \Theta (\mathrm{A}) \otimes \Theta (\mathrm{A}) \xrightarrow {\operatorname{Id} _ {\Theta (\mathrm{A})} \otimes c} \Theta (\mathrm{A}) \otimes \Theta (\mathrm{A}) \otimes \Theta (\mathrm{A}) \\ \Biggl \downarrow_ {j _ {2}} \\ (\mathrm{A} \otimes \mathrm{A}) ^ {*} \xrightarrow {^ t (\operatorname{Id} _ {\mathrm{A}} \otimes m _ {\mathrm{A}})} (\mathrm{A} \otimes \mathrm{A} \otimes \mathrm{A}) ^ {*}. \end{array}\tag{18}
\]

由于这个图以及 (15) 的交换性，对 \(f \in \Theta(\mathrm{A})\) 与 \(a, a', a'' \in A\)，我们有

\[
\langle j _ {3} \circ c ^ {\prime} (f), a \otimes a ^ {\prime} \otimes a ^ {\prime \prime} \rangle = \langle f, (a a ^ {\prime}) a ^ {\prime \prime} \rangle ;
\]

同理可证关系式

\[
\langle j _ {3} \circ c ^ {\prime \prime} (f), a \otimes a ^ {\prime} \otimes a ^ {\prime \prime} \rangle = \langle f, a (a ^ {\prime} a ^ {\prime \prime}) \rangle
\]

由于 A 中的乘法是结合的，我们有 \(j_{3} \circ c' = j_{3} \circ c''\)，因此 \(c' = c''\)，因为 \(j_{3}\) 是单射。

最后，公式 (16) 与 (17) 蕴含 \(\Theta(\mathrm{A})\) 以 \(\varepsilon\) 为余单位。

注 1. —— 设 \((\mathrm{V},\pi)\) 是代数 A 的一个有限维线性表示。取 V 的一组基 \((e_{1},\ldots,e_{n})\) 以及对偶基 \((e_{1}^{*},\ldots,e_{n}^{*})\)（\(V^{*}\) 的对偶基）。由引理 2 的证明，我们有关系式

\[
c (c _ {\pi} (x, x ^ {*})) = \sum_ {k = 1} ^ {n} c _ {\pi} (e _ {k}, x ^ {*}) \otimes c _ {\pi} (x, e _ {k} ^ {*})\tag{19}
\]

其中 \(x \in V\)，\(x^{*} \in V^{*}\)。对 \(1 \leqslant i, j \leqslant n\)，令 \(\pi_{ij} = c_{\pi}(e_j, e_i^*)\)。对每个 \(a \in A\)，\(\pi(a)\) 关于 V 的基 \((e_1, \ldots, e_n)\) 的矩阵等于 \((\pi_{ij}(a))\)。于是，对 \(1 \leqslant i \leqslant n\) 与 \(1 \leqslant j \leqslant n\)，我们有

\[
c (\pi_ {i j}) = \sum_ {k = 1} ^ {n} \pi_ {i k} \otimes \pi_ {k j}.\tag{20}
\]

定义 3. —— 设 C 是体 K 上的一个余代数，c 是它的余积。C 的余子代数是指 C 的这样的线性子空间 \(C_{1}\)：\(c(\mathrm{C}_{1})\) 包含于 \(C_{1} \otimes_{K} C_{1}\) 在 \(C \otimes_{K} C\) 中的典范像。

设 \(j\) 是 \(C_1\) 到 \(C\) 中的典范单射。按此定义，存在唯一的线性映射 \(c_1: C_1 \to C_1 \otimes_K C_1\)，使得我们有

\[
c \circ j = (j \otimes j) \circ c _ {1};\tag{21}
\]

这一关系式意味着 \(j\) 是从 \((\mathrm{C}_1, c_1)\) 到 \((\mathrm{C}, c)\) 的余代数态射（III, §11, No.~1, 第 574 页）。

若 C 是余结合的，则 \(C_{1}\) 也是余结合的。若 C 是余交换的，则 \(C_{1}\) 也是余交换的。若 C 有余单位 \(\varepsilon\)，则 \(\varepsilon\) 在 \(C_{1}\) 上的限制是 \(C_{1}\) 的一个余单位。

命题 9. —— 设 \(\Theta\) 是 \(\Theta(A)\) 的一个线性子空间。下列性质等价：

\begin{enumerate}
\def\labelenumi{(\roman{enumi})}
\item
  \(\Theta\) 是 \(\Theta(A)\) 的一个 (A, A)-子双模。
\item
  \(\Theta\) 是 \(\Theta(A)\) 的一个余子代数。
\item
  存在 A-模的类的一个遗传集 \(\mathcal{C}\)（这些 A-模都在 \(K\) 上有限维），使得 \(\Theta = \Theta_{\mathcal{C}}(A)\)。
\end{enumerate}

当这些性质成立时，(iii) 中所述的集合 C 是唯一确定的。

最后一个断言由 VIII, 第 388 页的推论得出：集合 C 由满足 \(\Theta_{\mathrm{E}}(\mathrm{A})\) 包含于 \(\Theta\) 的 K 上有限维 A-模 E 的类组成。

\begin{enumerate}
\def\labelenumi{(\roman{enumi})}
\setcounter{enumi}{2}
\item
  \(\Rightarrow\) (ii)：设 C 是 K 上有限维 A-模的类的一个遗传集。那么 \(\Theta_{\mathcal{C}}(\mathrm{A})\) 是有向族 \((\Theta_{\mathrm{E}}(\mathrm{A}))_{\mathrm{E}\in\mathcal{C}}\) 的并。由于 \(\Theta_{\mathrm{E}}(\mathrm{A})\) 是 \(\Theta(\mathrm{A})\) 的余子代数（对每个 \(E\in C\)）（VIII, 第 390 页，公式 (19)），故对 \(\Theta_{\mathcal{C}}(\mathrm{A})\) 亦然。
\item
  \(\Rightarrow\) (i)：设 \(f\in \Theta (\mathrm{A})\)。设 \(f_{1}^{\prime},\ldots ,f_{n}^{\prime},f_{1}^{\prime \prime},\ldots ,f_{n}^{\prime \prime}\) 是 \(\Theta (\mathrm{A})\) 中满足 \(c(f) = \sum f_i^\prime \otimes f_i''\) 的元素。对 A 中的 \(a,b\)，我们有 \(f(ab) = \sum f_i'(a)f_i''(b)\)，因此
\end{enumerate}

\[
b f = \sum_ {i = 1} ^ {n} f _ {i} ^ {\prime \prime} (b) f _ {i} ^ {\prime}, \quad f a = \sum_ {i = 1} ^ {n} f _ {i} ^ {\prime} (a) f _ {i} ^ {\prime \prime}
\]

（VIII, 第 375 页，公式 (3) 与 (4)）。因此，\(\Theta(A)\) 的余子代数是一个 \((A, A)\)-子双模。

\begin{enumerate}
\def\labelenumi{(\roman{enumi})}
\tightlist
\item
  \(\Rightarrow\) (iii)：设 \(\Theta\) 是 \(\Theta(\mathrm{A})\) 的一个 (A, A)-子双模；令 C 是满足 \(\Theta_{\mathrm{E}}(\mathrm{A})\) 包含于 \(\Theta\) 的 K 上有限维 A-模 E 的类所成的集合。集合 C 是遗传的（VIII, 第 377 页，注 2），且按构造有 \(\Theta_{\mathcal{C}}(\mathrm{A}) \subset \Theta\)。设 \(f \in \Theta\) 且 E = Af。那么 E 在 K 上有限维。于是，E 上的每个线性型都形如 \(u_{a}: g \mapsto g(a)\)，其中 a 属于 A（II, §7, No.~5, 第 302 页，推论 2）。而今，对 \(a, b, x \in A\)，我们有
\end{enumerate}

\[
c _ {\mathrm{E}} (a f, u _ {b}) (x) = \langle u _ {b}, x a f \rangle = f (b x a) = a f b (x),
\]

从而 \(c_{\mathrm{E}}(af, u_{b}) = afb\)。于是我们有 \(\Theta_{\mathrm{E}}(\mathrm{A}) \subset \Theta\)。因此 A-模 E 是 C 型的，且 f 是它的一个系数。于是我们有 \(\Theta \subset \Theta_{\mathcal{C}}(\mathrm{A})\)，最终得 \(\Theta = \Theta_{\mathcal{C}}(\mathrm{A})\)。

注 2. —— 设 \(\Theta\) 是 \(\Theta(A)\) 的一个余子代数。我们赋予 \(K\) 离散拓扑，并赋予代数 \(A\) 使映射 \(f: A \to K\)（其中 \(f\) 取遍 \(\Theta\)）都连续的最粗拓扑（Gen. Top., I, §2, No.~2, 第 175 页）。这个拓扑赋予 \(A\) 一个拓扑 \(K\)-模的结构。双边理想 \(\mathfrak{a}\)（\(A\) 的双边理想）是开的，当且仅当它具有有限余维数且其正交 \(\mathfrak{a}'\) 包含于 \(\Theta\)。由 VIII, 第 379 页的命题 3，\(A\) 的开双边理想构成 \(0\) 在 \(A\) 中的一个基本邻域系。因此 \(A\) 上的拓扑与其环结构相容。

设 C 是满足 \(\Theta = \Theta_{C}\) 的 K 上有限维 A-模的类的遗传集（VIII, 第 391 页，命题 9）。设 E 是 K 上有限维的左 A-模。我们赋予它离散拓扑。下列性质等价：

\begin{enumerate}
\def\labelenumi{(\roman{enumi})}
\item
  A-模 E 是 C 型的。
\item
  A-模 E 的零化理想在 A 中是开的。
\item
  从 A × E 到 E 的映射 \((a, x) \mapsto ax\) 是连续的。
\end{enumerate}

最后一个性质意味着 \(\mathrm{E}\) 是一个拓扑 A-模。

设 \(\Theta^{*}\) 是余代数 \(\Theta\) 的对偶代数。我们赋予 \(\Theta^{*}\) 使映射 \(\varphi \mapsto \varphi(u)\)（从 \(\Theta^{*}\) 到 K，其中 \(u\) 取遍 \(\Theta\)）都连续的最粗拓扑。代数 \(\Theta^{*}\) 上的拓扑与 \(\Theta^{*}\) 上的加群结构相容。\(\Theta^{*}\) 中形如 \(\Theta_{\mathrm{E}}(\mathrm{A})\)（其中 E 为 \(\mathcal{C}\) 型 A-模）的集合的正交构成 0 的一个基本邻域系。而今，这样的集合是 \(\Theta\) 的余子代数，故其正交是 \(\Theta^{*}\) 的理想。因此 \(\Theta^{*}\) 上的拓扑与其环结构相容。典范代数同态 \(\eta: A \to \Theta^{*}\)（它把 \(a \in A\) 映为线性型 \(f \mapsto f(a)\)，即 \(\Theta\) 上的线性型）定义了从 A 的完备豪斯多夫空间 \(\widehat{A}\)（Gen. Top., II, §3, No.~7, 第 191 页）到 \(\Theta^{*}\) 的一个同构。

\BourbakiExercises{习题}

\begin{enumerate}
\def\labelenumi{\arabic{enumi})}
\tightlist
\item
  设 \(\mathrm{K}\) 是一个交换域。设 \(\mathrm{A}\) 是一个 K-代数，它的基由单位元与两个元素 \(e, f\) 组成，满足 \(e^2 = e\)，\(ef = f\)，\(fe = f^2 = 0\)。a) 证明 \(\mathrm{A}\) 的正则表示与余正则表示不同构（计算 \(\operatorname{Tr}_{\mathrm{A / K}}(e)\) 与 \(\operatorname{Tr}_{\mathrm{A}^{\circ} / \mathrm{K}}(e)\)）。
\end{enumerate}

\begin{enumerate}
\def\labelenumi{\alph{enumi})}
\setcounter{enumi}{1}
\tightlist
\item
  证明正则表示包含一个子表示，它既不同构于余正则表示的任何子表示，也不同构于余正则表示的任何商表示。此外，存在 A 的一个单表示，它不同构于正则表示的任何子表示。
\end{enumerate}

\begin{enumerate}
\def\labelenumi{\arabic{enumi})}
\setcounter{enumi}{1}
\tightlist
\item
  设 \(A\) 是一个 K-代数。我们把对偶 \(A^*\)（\(A\) 的对偶）视为一个左 \(A\)-模。
\end{enumerate}

\begin{enumerate}
\def\labelenumi{\alph{enumi})}
\item
  映射 \(E \mapsto E'\) 是从 \(A^{*}\) 的 A-子模所成的集合到 A 的右理想所成的集合上的一个满射。当 E 是这样一个子模时，\(E'\) 在 \(A^{*}\) 中的正交是 \(A^{*}\) 的一个包含 E 的子模。
\item
  证明 \(\mathbf{A}^*\) 的底座在 A 中的正交是 A 的根。
\item
  证明在 \(A^{*}\) 中，\(A_{d}\) 的底座的正交是 A-模 \(A^{*}\) 的根。由此推出：\(A^{*}\) 是半单 A-模，当且仅当环 A 是半单的。
\end{enumerate}

¶ 3) 设 A 是一个阿廷环，\(\overline{\mathrm{A}}\) 是环 \(\mathrm{A} / \Re (\mathrm{A})\)。我们取 A 的一个分解，把它分解为不可分解理想 \(\mathrm{Ae}_i\)（\(i\in \mathrm{I}\)）的直和，其中诸 \(e_i\) 是正交幂等元（VIII, 第 179 页，习题 6, d)）。用 \(\mathrm{I} = \cup_{\alpha \in \mathrm{R}}\mathrm{I}_{\alpha}\) 表示 I 的这样一个分拆：I 的两个元素 \(i,j\) 属于同一个集合 \(\mathrm{I}_{\alpha}\)，当且仅当模 \(\mathrm{Ae}_i\) 与 \(\mathrm{Ae}_j\) 同构。\(\overline{\mathrm{A}}\)-模 \(\bigoplus_{i\in \mathrm{I}_{\alpha}}\overline{\mathrm{A}}\overline{e}_i\) 是 \(\overline{\mathrm{A}}\) 的同型分量（VIII, 第 179 页，习题 6, d) 与 7, a)）

\begin{enumerate}
\def\labelenumi{\alph{enumi})}
\item
  证明：若环 A 是自内射的（VIII, 第 53 页，习题 8），则存在 R 的一个置换 \(\sigma\)，使得每个右理想 \(e_i\mathrm{A}\)（对 \(i\in \mathrm{I}_{\alpha}\)）都包含唯一一个同构于 \(\overline{e}_j\overline{\mathrm{A}}\) 的极小右理想（对 \(j\in \mathrm{I}_{\sigma (\alpha)}\)）（注意幂等元 \(e\) 的右零化理想是 \((1 - e)\mathrm{A}\)）。由此推出 A 的左底座 \(\mathfrak{s}\) 包含于其右底座 \(\mathfrak{t}\)（证明每个右理想 \(e_i\mathfrak{s}\) 都是极小的，为此证明它是一个极大左理想的零化理想，参见 VIII, 第 178 页，习题 2）。最后断定 \(\mathfrak{s} = \mathfrak{t}\)，并且每个左理想 \(\mathrm{Ae}_i\)（对 \(i\in \mathrm{I}_{\sigma (\alpha)}\)）都包含唯一一个同构于 \(\overline{\mathrm{A}}\overline{e}_j\) 的极小左理想（对每个 \(j\in \mathrm{I}_{\alpha}\)）。
\item
  反之，若每个左理想 \(Ae_{i}\) 都包含唯一的极小左理想，且每个右理想 \(e_{i}A\) 都包含唯一的极小右理想，则环 A 是自内射的。（利用 VIII, 第 179 页的习题 7 证明 s 的每个足都是一个极小双边理想。由此推出 t 包含 s，进而 s = t。最后证明 VIII, 第 53 页习题 8 的条件 \(\beta\) ) 得到满足。）
\item
  设 A 是一个交换域上的有限次数代数，并且是一个自内射环。设 \(\alpha \in R\)，\(i \in I_{\alpha}\)，\(j \in I_{\sigma(\alpha)}\)。证明与 A-模 \(Ae_{j}\) 相伴的 A 的线性表示同构于 \(A^{\circ}\) 的线性表示（与右 A-模 \(e_{i}A\) 相伴的线性表示）的转置。（设 \(U_{i}\) 是 \(A^{*}\) 的作为 \((1 - e_{i})A\) 的正交的那个左子模。由 VIII, 第 179 页的习题 7, d) 推出 \(U_{i}\) 同构于 \(Ae_{j}\) 的一个商模，从而有 long \(e_{i}A \leqslant long e_{j}A\)。断定这两个长度相等，于是 \(U_{i}\) 同构于 \(Ae_{j}\)。利用 b) 与习题 2, a) 证明该命题的一个逆命题。）
\end{enumerate}

\begin{enumerate}
\def\labelenumi{\arabic{enumi})}
\setcounter{enumi}{3}
\tightlist
\item
  保留习题 3 的假设。
\end{enumerate}

\begin{enumerate}
\def\labelenumi{\alph{enumi})}
\item
  证明：若 A 是一个弗罗贝尼乌斯环（VIII, 第 53 页，习题 9），则有 \(\operatorname{Card}(\mathrm{I}_{\alpha}) = \operatorname{Card}(\mathrm{I}_{\sigma(\alpha)})\)（对每个 \(\alpha \in R\)）。并证明其逆。
\item
  一个交换域上的有限次数代数是弗罗贝尼乌斯环，当且仅当它的正则表示同构于它的余正则表示（利用习题 3, c)）。
\end{enumerate}

\begin{enumerate}
\def\labelenumi{\arabic{enumi})}
\setcounter{enumi}{4}
\tightlist
\item
  设 K 是一个交换域，A 是一个 K-代数。我们赋予 \(A^{*}\) 其典范左 A-模结构。对 \(t \in A^{*}\)，用 \(\varphi_{t}\) 表示映射 \((x, y) \mapsto t(xy)\)（从 \(A \times A\) 到 K 的映射）。
\end{enumerate}

\begin{enumerate}
\def\labelenumi{\alph{enumi})}
\item
  证明这定义了一个 K-线性同构，从 \(A^{*}\) 到由 K-双线性映射 \(\varphi: A \times A \to K\)（满足 \(\varphi(xy, z) = \varphi(x, yz)\)，对 \(x, y, z \in A\)）所成的集合。
\item
  我们说元素 \(t \in A^{*}\) 是可逆的，如果同态 \(a \mapsto at\)（从 A 到 \(A^{*}\)）是双射。设 A 是一个有限次数的 K-代数。证明下列条件等价：
\end{enumerate}

\begin{enumerate}
\def\labelenumi{(\roman{enumi})}
\item
  A 是弗罗贝尼乌斯环（习题 4）。
\item
  A 容许一个可逆线性型。
\item
  A 容许一个核不包含 \(\mathbf{A}\) 的任何左理想的线性型。
\end{enumerate}

\begin{enumerate}
\def\labelenumi{\arabic{enumi})}
\setcounter{enumi}{5}
\tightlist
\item
  设 K 是一个交换域。我们说 K-代数 A 是对称的，如果它容许一个可逆的迹型线性型（VIII, 第 115 页，习题 3）。
\end{enumerate}

\begin{enumerate}
\def\labelenumi{\alph{enumi})}
\item
  证明对称代数是有限次数的且是弗罗贝尼乌斯环。
\item
  证明 \(\mathrm{K}\) 上有限群的群代数是对称的。
\item
  \(\mathrm{K}\) 上有限次数的半单代数是对称的。
\end{enumerate}

\section{有限群的线性表示}

在本节中，G 是一个群，K 是一个交换环。若 G 是有限群，则用 \(|G|\) 表示 K 的元素 \(\operatorname{Card}(G) \cdot 1\)。从第 5 节起，我们假设群 G 是有限的，环 K 是一个代数闭域，且 \(|G|\) 不为零。

\subsection{线性表示}

定义 1. —— 设 M 是一个 K-模。G 在 M 中的线性表示是指从 G 到线性群 GL(M) 的一个群同态（II, §1, No.~2, 第 195 页）。

设 \(\pi\) 是 G 在 K-模 M 中的一个线性表示。我们也说偶 \((M,\pi)\) 是 G 的一个线性表示。当 K 非零且 M 是自由 K-模时，M 的维数称为表示 \(\pi\) 的次数（或维数）。

回顾一下，典范映射 \(g \mapsto e_{g}\)（从 G 到群 G 的代数 K{[}G{]} 的典范映射）是 K-模 K{[}G{]} 的一组基。滥用记号，我们也把元素 \(e_{g}\)（K{[}G{]} 的元素）记作 g。给定 G 的一个线性表示 \(\pi : G \to \mathbf{GL}(M)\)，存在唯一的从 K{[}G{]} 到 \(\operatorname{End}_{\mathrm{K}}(\mathrm{M})\) 的 K-代数同态延拓 \(\pi\)（III, §2, No.~6, 第 447 页，例）；我们也把这个延拓记作 \(\pi\)。它是代数 K{[}G{]} 在 M 中的一个表示（VIII, 第 373 页，定义 1），从而 M 是一个 K{[}G{]}-模。反之，每个线性表示 \(\rho\)（代数 K{[}G{]} 在 M 中的线性表示）都定义 G 在 M 中的一个线性表示，它由 \(g \mapsto \rho(g)\) 给出。

我们将不受拘束地把有关 K{[}G{]}-模结构的术语用于群 G 的线性表示。例如，我们说到子表示、单表示、表示的直和，等等。设 \((\mathrm{M},\pi)\) 与 \((\mathrm{M}^{\prime},\pi^{\prime})\) 是 G 的线性表示；用 \(\operatorname{Hom}_{\mathrm{G}}(\pi,\pi^{\prime})\) 表示 K-模 \(\operatorname{Hom}_{\mathrm{K[G]}}(\mathrm{M},\mathrm{M}^{\prime})\)（从 M 到 M' 的 K{[}G{]}-模同态所成的 K-模）。\(^{(1)}\)

映射 \(f\)（从 \(G\) 到 \(K\) 的映射）称为中心函数，如果 \(f(gg') = f(g'g)\) 对每一对 \((g, g')\)（\(G\) 中元素的每一对）成立；等价地，\(f(ghg^{-1}) = f(h)\) 对每一对元素 \(g, h\)（\(G\) 的元素对）成立。因此，中心函数就是 \(G\) 上在每个共轭类上取常值的函数。它们构成从 \(G\) 到 \(K\) 的映射空间的一个子模，记作 \(\mathcal{Z}_K(G)\)。当 \(G\) 有限时，K-代数 \(\mathcal{Z}_K(G)\) 是一个自由 K-模，其维数为 \(G\) 的共轭类的个数。中心 \(Z(K[G])\)（代数 \(K[G]\) 的中心）由元素 \(a = \sum a_g g\)（\(K[G]\) 中的元素）组成，这些元素满足 \(hah^{-1} = a\)（对每个 \(h \in G\)）。而今，我们有

\[
h a h ^ {- 1} = \sum_ {g \in \mathrm{G}} a _ {h ^ {- 1} g h} g;
\]

因此，当 G 有限时，中心 \(Z(K[G])\)（代数 K{[}G{]} 的中心）由中心函数组成。

设 \((\mathrm{M},\pi)\) 是 G 的一个线性表示。假设 M 是一个有限维自由 K-模。\(\pi\) 的迹是 K{[}G{]} 的表示（与 \(\pi\) 相伴的表示）的迹，即（VIII, 第 382 页）线性型 \(a\mapsto\mathrm{Tr}(\pi(a))\)（K{[}G{]} 上的线性型）。这个线性型由映射 \(g\mapsto\mathrm{Tr}(\pi(g))\)（从 G 到 K 的映射）确定，我们把该映射称为表示 \(\pi\) 的特征标，记作 \(\chi_{\pi}\)。表示的特征标是一个中心函数（II, §4, No.~3, 第 273 页，命题 3）。

设 \(M\) 与 \(M'\) 是有限维自由 \(K\)-模。设 \(\pi\) 与 \(\pi'\) 分别是 \(G\) 在 \(M\) 与 \(M'\) 中的线性表示。那么 \(M \oplus M'\) 是一个有限维自由 \(K\)-模，并且由 III, §9, No.~2, 第 543 页的命题 1，我们有

\[
\chi_ {\pi \oplus \pi^ {\prime}} = \chi_ {\pi} + \chi_ {\pi^ {\prime}}.
\]

更一般地，设

\[
0 \longrightarrow \mathrm{M} ^ {\prime} \longrightarrow \mathrm{M} \longrightarrow \mathrm{M} ^ {\prime \prime} \longrightarrow 0
\]

是 K{[}G{]}-模的一个正合列，并设 \(\pi\)、\(\pi'\)、\(\pi''\) 分别是与 M、\(M'\)、\(M''\) 相伴的 G 的线性表示。假设 K-模 \(M'\) 与 \(M''\) 都是自由且有限维的。那么 M 亦然（II, §1, No.~11, 第 218 页，命题 21），并且我们有

\[
\chi_ {\pi} = \chi_ {\pi^ {\prime}} + \chi_ {\pi^ {\prime \prime}}.\tag{1}
\]

命题 1. —— 假设 K 是一个交换域。设 \(\pi\) 与 \(\pi'\) 是 G 的有限维半单 K-线性表示。

\begin{enumerate}
\def\labelenumi{\alph{enumi})}
\item
  若 \(\pi(g)\) 与 \(\pi'(g)\) 的特征多项式对每个 \(g \in G\) 都相同，则表示 \(\pi\) 与 \(\pi'\) 同构。
\item
  再设 K 的特征为 0。若特征标 \(\chi_{\pi}\) 与 \(\chi_{\pi'}\) 相等，则表示 \(\pi\) 与 \(\pi'\) 同构。
\end{enumerate}

断言 a) 由 VIII, 第 386 页的推论 1 得出；断言 b) 由 VIII, 第 384 页的推论的 a) 得出。

例. —— 1) G 的单位表示或平凡表示是表示 \((\mathrm{K}, \varepsilon)\)，其中 \(\varepsilon(g) = \operatorname{Id}_{\mathrm{K}}\) 对每个 \(g \in G\) 成立。它的特征标是取常值 1 的常值函数。

\begin{enumerate}
\def\labelenumi{\arabic{enumi})}
\setcounter{enumi}{1}
\tightlist
\item
  G 的（左）正则表示是表示 \(\gamma\)（G 在 K{[}G{]} 中的表示），由 \(\gamma(g)(x)=gx\)（对 \(g\in G\) 与 \(x\in K[G]\)）定义。它对应于代数 K{[}G{]} 的左正则表示（VIII, 第 375 页）。设群 G 是有限的。此时正则表示是有限次数的。对 G 的每个不等于单位元的元素 g，g 在 K{[}G{]} 中的左乘关于典范基的矩阵是一个无不动点置换的矩阵。因此我们有
\end{enumerate}

\[
\chi_ {\boldsymbol {\gamma}} (g) = \left\{ \begin{array}{l l} | \mathrm{G} | & \text { if   g   is   the   identity   element }, \\ 0 & \text { otherwise }. \end{array} \right.\tag{2}
\]

我们类似地定义 G 的右正则表示。G 的双正则表示是表示 \(\rho\)（G × G 在 K{[}G{]} 中的表示），由 \(\rho(g, g')(x) = gxg'^{-1}\)（对 \((g, g') \in \mathrm{G} \times \mathrm{G}\) 与 \(x \in \mathrm{K}[\mathrm{G}]\)）定义。

\begin{enumerate}
\def\labelenumi{\arabic{enumi})}
\setcounter{enumi}{2}
\item
  给定 G 的一个线性表示 \((\mathrm{M},\pi)\)，逆步表示或对偶表示 \(\pi^{\vee}\)（\(\pi\) 的逆步表示或对偶表示）是 G 在与 M 对偶的 K-模中的表示，由关系式 \(\pi^{\vee}(g)={}^{t}\pi(g^{-1})\)（对每个 \(g\in G\)）定义（参见 II, §2, No.~5, 第 234 页）。若 M 是有限维自由 K-模，则它的对偶亦然，并且 \(\chi_{\pi^{\vee}}(g)=\chi_{\pi}(g^{-1})\) 对每个 \(g\in G\) 成立。
\item
  设 \((\mathrm{M},\pi)\) 与 \((\mathrm{M}^{\prime},\pi^{\prime})\) 是 G 的线性表示。在 VIII, 第 198 页的例 1 中，我们定义了一个 K{[}G{]}-模结构（在 \(M\otimes_{K}M^{\prime}\) 上）。相应的线性表示称为 \(\pi\) 与 \(\pi^{\prime}\) 的张量积，记作 \(\pi\otimes\pi^{\prime}\)。对 \(g\in G\)、\(x\in M\) 与 \(x^{\prime}\in M^{\prime}\)，有 \(\big((\pi\otimes\pi^{\prime})(g)\big)(x\otimes x^{\prime})=\pi(g)x\otimes\pi^{\prime}(g)x^{\prime}\)。
\end{enumerate}

若 M 与 M' 是有限维自由 K-模，则由 III, §9, No.~2, 第 543 页的命题 2，我们有

\[
\chi_ {\pi \otimes \pi^ {\prime}} = \chi_ {\pi} \chi_ {\pi^ {\prime}}.\tag{3}
\]

\begin{enumerate}
\def\labelenumi{\arabic{enumi})}
\setcounter{enumi}{4}
\tightlist
\item
  设 G 是两个群的乘积 \(G' \times G''\)。从 \(K[G'] \otimes_{K} K[G'']\) 到 K{[}G{]} 的、把 \(g' \otimes g''\) 映为 \((g', g'')\)（对 \(g' \in G'\) 与 \(g'' \in G''\)）的 K-线性映射是一个代数同构。设 \((M', \pi')\) 是 \(G'\) 的一个线性表示，\((M'', \pi'')\) 是 \(G''\) 的一个线性表示。\(\pi'\) 与 \(\pi''\) 的外张量积记作 \(\pi' \boxtimes \pi''\)，它是 \(G' \times G''\) 在向量空间 \(M' \otimes M''\) 中的表示，由 \((\pi' \boxtimes \pi'')(g', g'') = \pi'(g') \otimes \pi''(g'')\)（对 \((g', g'') \in G' \times G''\)）定义。若 \(M'\) 与 \(M''\) 是有限维自由 K-模，则 \(M' \otimes_{K} M''\) 是有限维自由 K-模，且表示 \(\pi' \boxtimes \pi''\) 的特征标由公式
\end{enumerate}

\[
\chi_ {\pi^ {\prime} \boxtimes \pi^ {\prime \prime}} (g ^ {\prime}, g ^ {\prime \prime}) = \chi_ {\pi^ {\prime}} (g ^ {\prime}) \chi_ {\pi^ {\prime \prime}} (g ^ {\prime \prime})
\]

（对 \(g' \in G'\) 与 \(g'' \in G''\)；III, §9, No.~2, 第 542 页，命题 2）。

\begin{enumerate}
\def\labelenumi{\arabic{enumi})}
\setcounter{enumi}{5}
\tightlist
\item
  设 \((V,\pi)\) 是 G 的一个线性表示，且 V 是有限维自由 K-模。我们用公式定义一个表示 \(\rho\)（\(G\times G\) 在 \(\operatorname{End}_{K}(V)\) 中的表示）：
\end{enumerate}

\[
\rho (g, g ^ {\prime}) (u) = \pi (g ^ {\prime}) \circ u \circ \pi (g ^ {- 1}).
\]

II, §4, No.~2, 第 271 页的 K-模同构 \(\theta_{\mathrm{V}}: \mathrm{V}^{*} \otimes_{\mathrm{K}} \mathrm{V} \to \mathrm{End}_{\mathrm{K}}(\mathrm{V})\) 是从外张量积 \(\pi^{\vee} \boxtimes \pi\) 到 \(\rho\) 的一个表示同构。映射 \(g \mapsto \rho(1, g)\)（相应地，\(g \mapsto \rho(g, 1)\)）是 G 的一个表示，它同构于 \(\pi^{\dim_{\mathrm{K}} \mathrm{V}}\)（相应地，\((\pi^{\vee})^{\dim_{\mathrm{K}} \mathrm{V}}\)）。

设 L 是一个交换 K-代数，\((\mathrm{M},\pi)\) 是群 G 的一个线性表示。群同态 \(\pi_{(\mathrm{L})}:\mathrm{G}\to\mathbf{GL}(\mathrm{M}_{(\mathrm{L})})\)（由 \(g\mapsto\mathrm{Id}_{\mathrm{L}}\otimes\pi(g)\) 定义）是 G 在 L-模 \(\mathrm{M}_{(\mathrm{L})}\) 中的一个线性表示，称为由表示 \(\pi\) 经标量环从 K 扩张到 L 而得到的 G 的线性表示。

设 K 是一个体，L 是一个非零交换 K-代数。设 \((M,\pi)\) 与 \((M',\pi')\) 是 G 的线性表示。表示 \(\pi\) 与 \(\pi'\) 同构，当且仅当 \(\pi_{(L)}\) 与 \(\pi'_{(L)}\) 同构（VIII, 第 37 页，定理 3）。

再设代数 L 是 K 的一个扩张。考虑 VIII, 第 198 页例 1 中定义的环 \(\mathrm{R}_{\mathrm{K}}(\mathrm{G})\) 与 \(\mathrm{R}_{\mathrm{L}}(\mathrm{G})\)。标量扩张定义了一个环同态

\[
u: \mathrm{R} _ {\mathrm{K}} (\mathrm{G}) \longrightarrow \mathrm{R} _ {\mathrm{L}} (\mathrm{G}).
\]

这个同态是单射，且元素 \(\xi \in \mathrm{R}_{\mathrm{K}}(\mathrm{G})\) 是有效的，当且仅当 \(u(\xi)\) 是有效的（VIII, 第 195 页，定理 1）。

\subsection{马施克定理}

定理 1. —— 设群 G 是有限的。设 M 是一个 K{[}G{]}-模，N 是 M 的一个 K{[}G{]}-子模。假设 N 是 K-模 M 的一个直因子，且 \textbar G\textbar{} 在 K 中可逆。那么 N 是 K{[}G{]}-模 M 的一个直因子。

设 p 是 M 中以 N 为像的一个 K-线性投影算子。我们如下定义 K-模 M 的一个自同态 q：

\[
q (x) = | \mathrm{G} | ^ {- 1} \sum_ {g \in \mathrm{G}} g p (g ^ {- 1} x)
\]

对每个 \(x \in M\) 成立。由于 N 在 G 的作用下稳定，且 p 在 N 上诱导恒同，我们看到 q 把 M 映入 N 且在 N 上诱导恒同。

于是 K-模 M 是 q 的像 N 与 q 的核的直和。\(g q(x) = q(gx)\)（对每个 \(x \in M\) 与 \(g \in G\)），故 q 的核是 M 的一个 K{[}G{]}-子模。这就证明了 N 是 K{[}G{]}-模 M 的一个直因子。

推论 1（马施克）. —— 设群 G 有限且 K 是一个交换域。代数 K{[}G{]} 是半单的，当且仅当域 K 的元素 \textbar G\textbar{} 不为零。

设 \(|\mathrm{G}| \neq 0\)。由定理 1，每个 K{[}G{]}-子模（\(\mathrm{K}[\mathrm{G}]_s\) 的子模）都是直因子。因此 K{[}G{]} 是半单的。

反之，设 \(|G|\) 为零，并用 \(\varepsilon\) 表示元素 \(\sum_{g\in G} g\)（K{[}G{]} 的中心的元素）。我们有 \(\varepsilon \neq 0\) 但 \(\varepsilon^{2} = |G|\varepsilon = 0\)，故 K{[}G{]} 不是半单的（VIII, 第 153 页，命题 3, a) 与 VIII, 第 157 页，注 1）。

推论 2. —— 设群 G 有限且 \(|G|\) 在 K 中可逆。K{[}G{]}-模的一个正合列分裂，当且仅当它作为 K-模的正合列分裂。

给定 K{[}G{]}-模的一个正合列

\[
0 \longrightarrow \mathrm{M} ^ {\prime} \stackrel {f} {\longrightarrow} \mathrm{M} \longrightarrow \mathrm{M} ^ {\prime \prime} \longrightarrow 0,
\]

只需对态射 f 的像应用定理 1。

推论 3. —— 设群 G 有限且 \(|G|\) 在 K 中可逆。一个 K{[}G{]}-模是投射的，当且仅当它作为 K-模是投射的。

设 P 是一个 K{[}G{]}-模。若 P 是自由 K{[}G{]}-模 M 的一个直因子，则它更是自由 K-模 M 的一个直因子。其逆由推论 2 得出，参见 II, §2, No.~2, 第 231 页，命题 4, d)。

推论 4. —— a) 设群 G 有限且 K 是一个特征为 0 的交换域。G 的两个有限维线性表示同构，当且仅当它们有相同的特征标。

\begin{enumerate}
\def\labelenumi{\alph{enumi})}
\setcounter{enumi}{1}
\tightlist
\item
  设 K 是一个特征为素数 p 的完满域，且群 G 有限、基数与 p 互素。G 的一个有限维线性表示的特征标为零，当且仅当该表示同构于 p 个两两同构的表示的直和。
\end{enumerate}

在该推论的假设下，G 的每个有限维线性表示都是半单的。于是该推论由 VIII, 第 384 页的推论得出。

推论 5. —— 设群 G 有限且 K 是一个满足 \(|G| \neq 0\) 的交换域。设 \(\pi\) 与 \(\pi'\) 是 G 的有限维线性表示。那么 \(\pi\) 与 \(\pi'\) 同构，当且仅当对 G 中的每个 g，自同态 \(\pi(g)\) 与 \(\pi'(g)\) 有相同的特征多项式。

这由 VIII, 第 386 页的推论 1 得出。

\subsection{诱导表示与余诱导表示}

设 H 是群 G 的一个子群。

若 \((V,\pi)\) 是 G 的一个线性表示，则 \(\pi\) 在 H 上的限制是 H 在 V 中的一个线性表示；我们把它记作 \(\operatorname{Res}_{\mathrm{H}}^{\mathrm{G}}(\pi)\)。K{[}H{]}-模（与 \(\operatorname{Res}_{\mathrm{H}}^{\mathrm{G}}(\pi)\) 相伴的那个模）就是由 K{[}G{]}-模 V 经标量限制得到的模（II, §1, No.~13, 第 221 页）。若 V 是有限维自由 K-模，则 \(\operatorname{Res}_{\mathrm{H}}^{\mathrm{G}}(\pi)\) 的特征标是 \(\pi\) 的特征标在 H 上的限制。

设 \((\mathrm{M},\sigma)\) 是 H 的一个表示。

我们把 K{[}G{]} 视为一个 (K{[}G{]}, K{[}H{]})-双模，把 M 视为一个 K{[}H{]}-模。K{[}G{]}-模 \(\mathcal{T}(M) = K[G] \otimes_{K[H]} M\)（VIII, 第 58 页）定义了 G 的一个线性表示，记作 \(\operatorname{Ind}_{\mathrm{H}}^{\mathrm{G}}(\sigma)\)，称为由 \(\sigma\) 诱导的 G 的表示。若 \((\mathrm{V}, \pi)\) 是 G 的一个线性表示，则 K{[}H{]}-模 \(\mathcal{H}(\mathrm{V}) = \mathrm{Hom}_{\mathrm{K[G]}}(\mathrm{K[G]},\mathrm{V})\) 可与一个 K{[}H{]}-模等同，该模对应于表示 \(\mathrm{Res}_{\mathrm{H}}^{\mathrm{G}}(\pi)\)。于是，伴随态射（VIII, 第 59 页）给出一个从 \(\mathrm{Hom}_{\mathrm{H}}(\sigma ,\mathrm{Res}_{\mathrm{H}}^{\mathrm{G}}(\pi))\) 到 \(\mathrm{Hom}_{\mathrm{G}}(\mathrm{Ind}_{\mathrm{H}}^{\mathrm{G}}(\sigma),\pi)\) 的 K-模同构，称为典范同构（"弗罗贝尼乌斯互反律"）。

我们把 K{[}G{]} 视为一个 (K{[}H{]}, K{[}G{]})-双模。K{[}G{]}-模 \(\mathcal{H}(M) = \operatorname{Hom}_{\mathrm{K}[\mathrm{H}]}(\mathrm{K}[\mathrm{G}], \mathrm{M})\) 定义了 G 的一个表示，记作 \(\operatorname{Coind}_{\mathrm{H}}^{\mathrm{G}}(\sigma)\)，称为由 \(\sigma\) 余诱导的 G 的表示。若 \((\mathrm{V}, \pi)\) 是 G 的一个线性表示，则 K{[}H{]}-模 \(\mathcal{T}(\mathrm{V}) = \mathrm{K}[\mathrm{G}] \otimes_{\mathrm{K}[\mathrm{G}]} \mathrm{V}\) 可与一个 K{[}H{]}-模等同，该模对应于表示 \(\operatorname{Res}_{\mathrm{H}}^{\mathrm{G}}(\pi)\)。于是，伴随态射（同前引文）给出一个从 \(\operatorname{Hom}_{\mathrm{H}}(\operatorname{Res}_{\mathrm{H}}^{\mathrm{G}}(\pi), \sigma)\) 到 \(\operatorname{Hom}_{\mathrm{G}}(\pi, \operatorname{Coind}_{\mathrm{H}}^{\mathrm{G}}(\sigma))\) 的 K-模同构，称为典范同构。

设 \(\varepsilon : \mathrm{K}[\mathrm{G}] \to \mathrm{K}[\mathrm{H}]\) 是由关系式 \(\varepsilon(h) = h\)（若 \(h \in \mathrm{H}\)）与 \(\varepsilon(g) = 0\)（若 \(g \in \mathrm{G} - \mathrm{H}\)）刻画的 K-模同态。映射 \(\varepsilon\) 是一个 \((\mathrm{K}[\mathrm{H}], \mathrm{K}[\mathrm{H}])\)-双模同态。设 \((\mathrm{M}, \sigma)\) 是 H 的一个线性表示。映射 \(v \mapsto v \circ \varepsilon\)（从 \(\mathrm{Hom}_{\mathrm{K}[\mathrm{H}]}(\mathrm{K}[\mathrm{H}], \mathrm{M})\) 到 \(\mathrm{Hom}_{\mathrm{K}[\mathrm{H}]}(\mathrm{K}[\mathrm{G}], \mathrm{M})\)）是一个 \(\mathrm{K}[\mathrm{H}]\)-模同态。把 M 与 \(\mathrm{Hom}_{\mathrm{K}[\mathrm{H}]}(\mathrm{K}[\mathrm{H}], \mathrm{M})\) 等同，我们就得到一个 \(\mathrm{K}[\mathrm{H}]\)-模同态，从 M 到 \(\mathrm{Res}_{\mathrm{H}}^{\mathrm{G}}(\mathrm{Coind}_{\mathrm{H}}^{\mathrm{G}}(\sigma))\)。弗罗贝尼乌斯互反律把它变为一个模同态 \(\iota\)（\(\mathrm{K}[\mathrm{G}]\)-模的同态），从 \(\mathrm{Ind}_{\mathrm{H}}^{\mathrm{G}}(\sigma)\) 到 \(\mathrm{Coind}_{\mathrm{H}}^{\mathrm{G}}(\sigma)\)，它由关系式

\[
\iota (g \otimes m) (g ^ {\prime}) = \varepsilon (g ^ {\prime} g) m
\]

（对 \(g, g' \in G\) 与 \(m \in M\)）刻画。这个同态称为典范同态。

命题 2. —— 设 H 是 G 的一个子群，\((\mathrm{M},\sigma)\) 是 H 的一个线性表示。典范同态 \(\iota:\operatorname{Ind}_{\mathrm{H}}^{\mathrm{G}}(\sigma)\to\operatorname{Coind}_{\mathrm{H}}^{\mathrm{G}}(\sigma)\) 是单射。若子群 H 在 G 中有有限指数，则 \(\iota\) 是一个 K{[}G{]}-模同构。

设 \(S \subset G\) 是 \(G/H\) 的一个代表系。族 \((s)_{s \in S}\) 是右 \(K[H]\)-模 \(K[G]\) 的一组基。由此得出，映射 \(M^{(S)} \to \operatorname{Ind}_{H}^{G}(\sigma)\)（由 \((m_s)_{s \in S} \mapsto \sum_{s \in S} s \otimes m_s\) 定义）是一个 \(K\)-模同构。对一切 \(s, s' \in S\) 与每个 \(m \in M\)，我们有关系式

\[
\iota (s ^ {\prime} \otimes m) (s ^ {- 1}) = \left\{ \begin{array}{l} m \text {   if   } s = s ^ {\prime}, \\ 0 \text {   otherwise }. \end{array} \right.
\]

由此得出 \(\iota'\) 是单射。

设 H 在 G 中有有限指数。此时集合 S 是有限的；设 \(\rho\) 是从 \(\operatorname{Coind}_{\mathrm{H}}^{\mathrm{G}}(\sigma)\) 到 \(\operatorname{Ind}_{\mathrm{H}}^{\mathrm{G}}(\sigma)\) 的映射，由 \(u \mapsto \sum_{s \in S} s \otimes u(s^{-1})\) 给出。它满足 \((\iota \circ \rho(u))(s^{-1}) = u(s^{-1})\)（对 \(u \in \operatorname{Coind}_{\mathrm{H}}^{\mathrm{G}}(\sigma)\) 与 \(s \in S\)）。由于族 \((s^{-1})_{s\in\mathrm{S}}\) 是左 K{[}H{]}-模 K{[}G{]} 的一组基，映射 \(\iota\circ\rho\) 是恒同映射。于是 \(\iota\) 是双射，其逆双射为 \(\rho\)。

设 H 是 G 的一个有限指数子群。设 u 是 H 上的一个中心函数；用 \(u^{0}\) 表示 G 上延拓 u 且在 G − H 的每一点取零的函数。设 S 是 G/H 的一个代表系。对 \(g \in G\)，令

\[
\operatorname{Ind} _ {\mathrm{H}} ^ {\mathrm{G}} (u) (g) = \sum_ {s \in \mathrm{S}} u ^ {0} (s ^ {- 1} g s).\tag{4}
\]

注意，对一切 \(x, g \in G\) 与每个 \(h \in H\)，有 \(u^{0}((xh)^{-1}gxh) = u^{0}(x^{-1}gx)\)。由此得出，\(\operatorname{Ind}_{\mathrm{H}}^{\mathrm{G}}(u)\) 是 G 上的一个不依赖于 S 的选取的中心函数。这样我们就定义了一个 K-线性映射 \(Ind_{H}^{G}\)，从 \(\mathcal{Z}_{\mathrm{K}}(\mathrm{H})\) 到 \(\mathcal{Z}_{\mathrm{K}}(\mathrm{G})\)。

命题 3. —— 设 H 是 G 的一个有限指数子群。设 \((M,\sigma)\) 是 H 的一个线性表示。假设 M 是一个有限维自由 K-模。用 \((V,\pi)\) 表示由 \(\sigma\) 诱导的 G 的表示。那么 K-模 V 是自由且有限维的，并且

\[
\chi_ {\pi} = \mathrm{Ind} _ {\mathrm{H}} ^ {\mathrm{G}} (\chi_ {\sigma}).\tag{5}
\]

设 S 是 G/H 的一个代表系。如上所见，线性映射 \(\mathrm{M}^{\mathrm{S}}\to\mathrm{Ind}_{\mathrm{H}}^{\mathrm{G}}(\mathrm{M})\)（由 \((m_{s})_{s\in\mathrm{S}}\mapsto\sum_{s\in\mathrm{S}}s\otimes m_{s}\) 给出）是一个 K-模同构。特别地，V 是一个有限维自由 K-模。对任意 \(g\in G\)，用 \(M_{g}\) 表示 M 在映射 \(m\mapsto g\otimes m\) 下的像。对一切 \(g,g'\in G\)，\(M_{g}=M_{g'}\) 当且仅当 \(gH=g'H\)，且 V 是诸子模 \(M_{s}\)（\(s\in S\)）的直和。对一切 \(g,g'\in G\)，我们还有 \(\pi(g)\mathrm{M}_{g'}=\mathrm{M}_{gg'}\)。对任意 \(g\in G\)，用 \(S_{g}\) 表示 S 中满足 \(s^{-1}gs\in H\) 的元素 s 所成的集合。对一切 \(g,g'\in G\)（满足 \(g'^{-1}gg'\in H\)），V 的自同构 \(\pi(g)\) 诱导一个自同构 \(\pi(g)_{g'}\)（\(M_{g'}\) 的自同构），并且我们有

\[
\operatorname{Tr} (\pi (g)) = \sum_ {s \in \mathrm{S} _ {g}} \operatorname{Tr} (\pi (g) _ {s}) = \sum_ {s \in \mathrm{S} _ {g}} \operatorname{Tr} (\sigma (s ^ {- 1} g s)).
\]

命题的最后一个断言由此得出。

\subsection{表示与格罗滕迪克群}

在本小节中，我们假设 K 是一个交换域。给定 G 的一个有限维线性表示 \((M,\pi)\)，用 \([\pi]\) 表示 K{[}G{]}-模 M 在格罗滕迪克环 \(\mathrm{R}_{\mathrm{K}}(\mathrm{G})\) 中的类（VIII, 第 198 页，例 1）。

由环 \(\mathrm{R}_{\mathrm{K}}(\mathrm{G})\) 上合成律的定义，若 \(\pi\) 与 \(\pi'\) 是 G 的有限维线性表示，则我们有

\[
[ \pi \oplus \pi^ {\prime} ] = [ \pi ] + [ \pi^ {\prime} ], \quad [ \pi \otimes \pi^ {\prime} ] = [ \pi ] [ \pi^ {\prime} ].
\]

环 \(R_{K}(G)\) 的单位元是 G 的单位表示的类（VIII, 第 399 页，例 1）。

环 \(R_{K}(G)\) 是一个自由 Z-模，以 K 上有限维单 K{[}G{]}-模的类所成的集合为基。当 \(|G|\) 在 K 中可逆时，两个有限维表示同构，当且仅当它们在 \(R_{K}(G)\) 中有相同的类（VIII, 第 190 页，推论，及第 401 页，推论 1）。

对 G 的每个有限维线性表示 \((M,\pi)\)，逆步表示 \(\pi^{\vee}\) 也是有限维的。表示 \(\pi\) 与 \((\pi^{\vee})^{\vee}\) 同构。若 \(\pi'\) 也是 G 的一个有限维线性表示，则表示 \((\pi\otimes\pi')^{\vee}\) 与 \(\pi^{\vee}\otimes\pi'^{\vee}\) 同构。若

\[
0 \longrightarrow \mathrm{M} ^ {\prime} \longrightarrow \mathrm{M} \longrightarrow \mathrm{M} ^ {\prime \prime} \longrightarrow 0
\]

是 K 上有限维 K{[}G{]}-模的一个正合列，则逆步表示给出 K{[}G{]}-模的一个正合列

\[
0 \longrightarrow \mathrm{M} ^ {\prime \prime \vee} \longrightarrow \mathrm{M} ^ {\vee} \longrightarrow \mathrm{M} ^ {\prime \vee} \longrightarrow 0
\]

（II, §7, No.~5, 第 299 页，定理 5 的推论）。由格罗滕迪克群的泛性质（VIII, 第 186 页，命题 4），存在一个自同构 \(c \mapsto c^{\vee}\)（环 \(\mathrm{R}_{\mathrm{K}}(\mathrm{G})\) 的自同构），它由 \([\pi]^{\vee} = [\pi^{\vee}]\)（对每个有限维线性表示 \(\pi\)，\(\mathrm{G}\) 的表示）刻画；我们有 \((c^{\vee})^{\vee} = c\)（对每个元素 \(c\)，\(\mathrm{R}_{\mathrm{K}}(\mathrm{G})\) 的元素）。

设 H 是 G 的一个子群。限制保持正合列。于是由格罗滕迪克群的泛性质（同前引文），存在一个群同态，记作 \(Res_{H}^{G}\)，它从 \(R_{K}(G)\) 到 \(R_{K}(H)\)，由关系式

\[
\mathrm{Res} _ {\mathrm{H}} ^ {\mathrm{G}} [ \pi ] = [ \mathrm{Res} _ {\mathrm{H}} ^ {\mathrm{G}} (\pi) ]\tag{6}
\]

（对 G 的每个有限维线性表示 \(\pi\)）刻画；它是一个环同态。若 H 在 G 中有有限指数，则由 II, §7, No.~7, 第 306 页的命题 14，我们可以类似地定义一个群同态 \(\mathrm{Ind}_{\mathrm{H}}^{\mathrm{G}}\)，从 \(\mathrm{R}_{\mathrm{K}}(\mathrm{H})\) 到 \(\mathrm{R}_{\mathrm{K}}(\mathrm{G})\)，它由关系式

\[
\mathrm{Ind} _ {\mathrm{H}} ^ {\mathrm{G}} [ \sigma ] = [ \mathrm{Ind} _ {\mathrm{H}} ^ {\mathrm{G}} (\sigma) ]\tag{7}
\]

对 H 的每个有限维线性表示 \(\sigma\) 成立。

由 VIII, 第 399 页的关系式 (1)、VIII, 第 400 页的关系式 (3) 以及上面提到的格罗滕迪克群的泛性质，存在一个环同态 \(\Theta_{\mathrm{G}}\)，从 \(\mathrm{R}_{\mathrm{K}}(\mathrm{G})\) 到中心函数的代数 \(\mathcal{Z}_{\mathrm{K}}(\mathrm{G})\)（\(\mathrm{G}\) 上中心函数的代数），它由 \(\Theta_{\mathrm{G}}([\pi]) = \chi_{\pi}\)（对每个有限维表示 \(\pi\)，\(\mathrm{G}\) 的表示）刻画。

若 H 是 G 的子群，则与群 G 和 H 相伴的同态 \(\Theta_{G}\) 与 \(\Theta_{H}\) 同运算 \(Res_{H}^{G}\) 与 \(Ind_{H}^{G}\) 相容（VIII, 第 402 页，以及 VIII, 第 404 页，命题 3）。

设 G 是两个群的乘积 \(G^{\prime} \times G^{\prime\prime}\)。由 VIII, 第 213 页注 2，存在一个 Z-线性映射 \(\kappa\)，从 \(\mathrm{R}_{\mathrm{K}}(\mathrm{G}^{\prime}) \otimes_{\mathbf{Z}} \mathrm{R}_{\mathrm{K}}(\mathrm{G}^{\prime\prime})\) 到群 \(\mathrm{R}_{\mathrm{K}}(\mathrm{G}^{\prime} \times \mathrm{G}^{\prime\prime})\)，它由关系式 \(\kappa([\pi^{\prime}] \otimes [\pi^{\prime\prime}]) = [\pi^{\prime} \boxtimes \pi^{\prime\prime}]\) 刻画，其中 \(\pi^{\prime}\)（相应地，\(\pi^{\prime\prime}\)）是 \(G^{\prime}\)（相应地，\(G^{\prime\prime}\)）的一个有限维表示，而 \(\pi^{\prime} \boxtimes \pi^{\prime\prime}\) 是外张量积（VIII, 第 400 页，例 5）。它是一个环同态。若域 K 是代数闭的，则映射 \(\kappa\) 是一个同构（VIII, 第 213 页，注 2）。

设 G 是两个群的乘积 \(G' \times G''\)。用 \(\psi\) 表示从 \(\mathcal{Z}_{\mathrm{K}}(\mathrm{G}') \otimes_{\mathrm{K}} \mathcal{Z}_{\mathrm{K}}(\mathrm{G}'' )\) 到 \(\mathcal{Z}_{\mathrm{K}}(\mathrm{G})\) 的、把 \(f' \otimes f''\) 映为函数 \((g', g'') \mapsto f'(g')f''(g'')\) 的同态。下图交换：

\[
\mathrm{R} _ {\mathrm{K}} (\mathrm{G} ^ {\prime}) \otimes_ {\mathbf {Z}} \mathrm{R} _ {\mathrm{K}} (\mathrm{G} ^ {\prime \prime}) \xrightarrow {\kappa} \mathrm{R} _ {\mathrm{K}} (\mathrm{G})
\]

\[
\bigg \downarrow \Theta_ {\mathrm{G} ^ {\prime}} \otimes \Theta_ {\mathrm{G} ^ {\prime \prime}} \qquad \qquad \bigg \downarrow \Theta_ {\mathrm{G}}
\]

\[
\mathcal {Z} _ {\mathrm{K}} (\mathrm{G} ^ {\prime}) \otimes_ {\mathrm{K}} \mathcal {Z} _ {\mathrm{K}} (\mathrm{G} ^ {\prime \prime}) \xrightarrow {\psi} \mathcal {Z} _ {\mathrm{K}} (\mathrm{G}).
\]

\subsection{傅里叶反演公式}

在本节余下部分，我们假设群 G 有限，且 K 是一个代数闭域，其特征不整除 G 的阶，从而 K 的元素 \(|G|\) 不为零。

代数 K{[}G{]} 是半单的（马施克定理）且有限维。用 \(\widehat{G}\) 表示单 K{[}G{]}-模的类所成的集合。对每个 \(\lambda\in\widehat{G}\)，选取 G 的一个线性表示 \((\mathrm{V}_{\lambda},\pi_{\lambda})\)，使相伴的 K{[}G{]}-模的类为 \(\lambda\)。集合 \(\widehat{G}\) 是有限的，且向量空间 \(V_{\lambda}\) 都是有限维的（VIII, 第 141 页，例）。对任意 \(\lambda\in\widehat{G}\)，用 \(d_{\lambda}\) 表示表示 \(\pi_{\lambda}\) 的次数，即 K-向量空间 \(V_{\lambda}\) 的维数，并用 \(\chi_{\lambda}\) 表示它的特征标。

用 \(\mathrm{F}(\widehat{\mathrm{G}})\) 表示乘积代数 \(\prod_{\lambda\in\widehat{G}}\operatorname{End}_{\mathrm{K}}(\mathrm{V}_{\lambda})\)，用 \(\overline{F}\) 表示从 K{[}G{]} 到 \(\mathrm{F}(\widehat{\mathrm{G}})\) 的映射，它由 \(\overline{\mathcal{F}}(a)=(\pi_{\lambda}(a))_{\lambda\in\widehat{G}}\) 定义。由于域 K 是代数闭的，映射 \(\overline{F}\) 是一个代数同构（同前引文）。

对每个 \(\lambda \in \widehat{\mathbf{G}}\)，代数 \(\mathrm{End}_{\mathrm{K}}(\mathrm{V}_{\lambda})\) 的维数是 \(d_{\lambda}^{2}\)；代数 \(\mathrm{K}[\mathrm{G}]\) 的维数是 \(\mathrm{Card}(\mathrm{G})\)。因此我们有关系式

\[
\mathrm{Card(G)} = \sum_ {\lambda \in \widehat {\mathrm{G}}} d _ {\lambda} ^ {2}.\tag{8}
\]

用 \(\tau\) 表示代数 K{[}G{]} 中的迹；按定义，迹 \(\tau(a)\)（元素 \(a\) 的迹，此元素属于 K{[}G{]}）是自同态 \(x \mapsto ax\)（K{[}G{]} 的自同态）的迹（III, §7, No.~7, 第 515 页）。设 \(a = \sum_{g \in G} a_g g\) 是 K{[}G{]} 的一个元素；由公式 (2)，我们有 \(\tau(ag^{-1}) = |G| a_g\)（对每个 \(g \in G\)），从而有关系式

\[
a = | \mathrm{G} | ^ {- 1} \sum_ {g \in \mathrm{G}} \tau (a g ^ {- 1}) g.\tag{9}
\]

用 \(\widehat{\tau}\) 表示代数 \(F(\widehat{G})\) 中的迹。设 \(A = (A_{\lambda})_{\lambda \in \widehat{G}}\) 是 \(F(\widehat{G})\) 的一个元素；我们有（参见 III, §9, No.~3, 第 545 页，例 3）

\[
\widehat {\tau} (\mathrm{A}) = \sum_ {\lambda \in \widehat {\mathrm{G}}} d _ {\lambda} \operatorname{Tr} (\mathrm{A} _ {\lambda}).\tag{10}
\]

由于映射 \(\overline{F}\) 是一个 K-代数同构，我们有 \(\widehat{\tau}\circ\overline{F}=\tau\)，因此

\[
\tau (a) = \widehat {\tau} (\overline {{\mathcal {F}}} (a)) = \widehat {\tau} \big ((\pi_ {\lambda} (a)) _ {\lambda \in \widehat {\mathrm{G}}} \big) = \sum_ {\lambda \in \widehat {\mathrm{G}}} d _ {\lambda} \operatorname{Tr} (\pi_ {\lambda} (a)) = \sum_ {\lambda \in \widehat {\mathrm{G}}} d _ {\lambda} \operatorname{Tr} _ {\lambda} (a)
\]

对每个 \(a\in \mathrm{K}[\mathrm{G}]\) 成立，也就是

\[
\tau = \sum_ {\lambda \in \widehat {\mathrm{G}}} d _ {\lambda} \operatorname{Tr} _ {\lambda}.\tag{11}
\]

于是，由 VIII, 第 399 页的 (2)，对 \(g \in \mathbf{G}\)，我们有

\[
\sum_ {\lambda \in \widehat {G}} d _ {\lambda}   \chi_ {\lambda} (g) = \left\{ \begin{array}{l l} | G | & \text { if   } g \text {   is   the   identity   element }, \\ 0 & \text { otherwise }. \end{array} \right.\tag{12}
\]

对 \(a \in \mathrm{K}[\mathrm{G}]\)，关系式 (9) 取如下形式：

\[
a = | \mathrm{G} | ^ {- 1} \sum_ {g \in \mathrm{G}} \sum_ {\lambda \in \widehat {\mathrm{G}}} d _ {\lambda} \left. \operatorname{Tr} (\pi_ {\lambda} (a) \pi_ {\lambda} (g ^ {- 1})) g; \right.\tag{13}
\]

由此得出，对元素 \(A = (A_{\lambda})_{\lambda \in \widehat{G}}\)（\(F(\widehat{G})\) 的每个元素），我们有

\[
\overline {{\mathcal {F}}} ^ {- 1} (\mathrm{A}) = | \mathrm{G} | ^ {- 1} \sum_ {g \in \mathrm{G}} \sum_ {\lambda \in \widehat {\mathrm{G}}} d _ {\lambda} \operatorname{Tr} (\mathrm{A} _ {\lambda} \pi_ {\lambda} (g ^ {- 1})) g\tag{14}
\]

（"傅里叶反演公式"）。

对 \(\mu \in \widehat{\mathrm{G}}\)，用 \(j_{\mu}:\mathrm{End}_{\mathrm{K}}(\mathrm{V}_{\mu})\longrightarrow \prod_{\lambda \in \widehat{\mathrm{G}}} \mathrm{End}_{\mathrm{K}}(\mathrm{V}_{\lambda})\) 表示满足 \(j_{\mu}(u) = (v_{\lambda})\) 的映射，其中 \(v_{\lambda} = 0\)（若 \(\lambda \neq \mu\)），而 \(v_{\mu} = u\)。由公式 (14)，我们有

\[
\overline {{\mathcal {F}}} ^ {- 1} (j _ {\mu} (u)) = | \mathrm{G} | ^ {- 1} d _ {\mu} \sum_ {g \in \mathrm{G}} \mathrm{Tr} (u \pi_ {\mu} (g ^ {- 1})) g.\tag{15}
\]

代数 \(\mathrm{F}(\widehat{\mathrm{G}})=\prod_{\lambda\in\widehat{\mathrm{G}}} \mathrm{End}_{\mathrm{K}}(\mathrm{V}_{\lambda})\) 的中心由诸族 \((a_{\lambda}1_{\mathrm{V}_{\lambda}})_{\lambda\in\widehat{\mathrm{G}}}\) 组成，其中 \((a_{\lambda})\) 是 K 的元素的一个族。它就是 \(\overline{{\mathcal{F}}}\) 作用于代数 K{[}G{]} 的中心所得到的像。于是该中心有一组基 \((e_{\lambda})_{\lambda\in\widehat{\mathrm{G}}}\)，它由关系式

\[
\pi_ {\lambda} (e _ {\mu}) = \delta_ {\lambda \mu} 1 _ {V _ {\lambda}}\tag{16}
\]

（对 \(\lambda, \mu \in \widehat{\mathbf{G}}\)）刻画，其中 \(\delta_{\lambda \mu}\) 是克罗内克符号。由公式 (15)，对每个 \(\mu \in \widehat{\mathbf{G}}\)，我们有

\[
e _ {\mu} = | \mathrm{G} | ^ {- 1} d _ {\mu} \sum_ {g \in \mathrm{G}} \chi_ {\mu} (g ^ {- 1}) g.\tag{17}
\]

这些元素满足关系式

\[
\sum_ {\lambda \in \widehat {\mathsf {G}}} e _ {\lambda} = 1, \quad e _ {\mu} ^ {2} = e _ {\mu}, \quad \mathrm{and} \quad e _ {\mu} e _ {\nu} = 0\tag{18}
\]

对一切 \(\mu, \nu \in \widehat{\mathrm{G}}\)（满足 \(\mu \neq \nu\)）；它们是 \(\mathrm{Z}(\mathrm{K}[\mathrm{G}])\) 的不可分解幂等元（VIII, 第 147 页，命题 15）

注. —— 设 \((V, \pi)\) 是 G 的一个线性表示。由马施克定理，K{[}G{]}-模 V 是半单的。对任意 \(\lambda \in \widehat{G}\)，用 \(V^{\lambda}\) 表示 \(\lambda\) 型同型分量（即 K{[}G{]}-模 V 的同型分量）；我们有 \(V = \bigoplus_{\lambda \in \widehat{G}} V^{\lambda}\)。由 VIII, 第 147 页的命题 15 与公式 (17)，与 V 的这个分解相伴的、以 \(V^{\lambda}\) 为像的 V 的投影算子等于

\[
\pi (e _ {\lambda}) = | \mathrm{G} | ^ {- 1} d _ {\lambda} \sum_ {g \in \mathrm{G}} \chi_ {\lambda} (g ^ {- 1}) \pi (g).\tag{19}
\]

把这个公式应用于 \(\pi = \pi_{\lambda}\)，我们得到 K 的元素 \(d_{\lambda} \cdot 1\) 不为零。我们将在后面（VIII, 第 420 页，推论 2）看到 \(d_{\mu}\) 整除 G 的基数。

设 \(\lambda\) 是 \(\widehat{\mathbf{G}}\) 的一个元素。映射 \(\pi_{\lambda}\)（从 K{[}G{]} 到 \(\mathrm{End}_{\mathrm{K}}(\mathrm{V}_{\lambda})\)）诱导一个 (K{[}G{]}, K{[}G{]})-双模同构，从 \(e_{\lambda}\mathrm{K}[\mathrm{G}]\) 到 \(\mathrm{End}_{\mathrm{K}}(\mathrm{V}_{\lambda})\)。由 VIII, 第 147 页的命题 15，K{[}G{]}-子模 \(e_{\lambda}\mathrm{K}[\mathrm{G}]\) 是 K{[}G{]} 的 \(\lambda\) 型同型分量。它也是 G 的右正则表示的 \(\lambda^{\vee}\) 型同型分量（VIII, 第 143 页，命题 11），同时又是双正则表示的 \(\lambda \times \lambda^{\vee}\) 型同型分量（VIII, 第 381 页，命题 4）。

\subsection{舒尔正交关系}

我们保留上一小节的记号。设 \(\lambda\) 是 \(\widehat{G}\) 的一个元素，\(u\) 与 \(v\) 是 \(\mathrm{End}_{\mathrm{K}}(\mathrm{V}_{\lambda})\) 的元素；由公式 (10)，我们有

\[
\hat {\tau} (j _ {\lambda} (u) j _ {\lambda} (v)) = d _ {\lambda} \operatorname{Tr} (u v).
\]

由于 \(\tau = \hat{\tau} \circ \overline{\mathcal{F}}\)，我们从公式 (2) 与 (15) 导出关系式

\[
d _ {\lambda} ^ {2} | \mathrm{G} | ^ {- 1} \sum_ {g \in \mathrm{G}} \operatorname{Tr} (u \pi_ {\lambda} (g)) \operatorname{Tr} (v \pi_ {\lambda} (g ^ {- 1})) = d _ {\lambda} \operatorname{Tr} (u v).
\]

我们先前已指出，在域 K 中 \(d_{\lambda}.1 \neq 0\)；由此得出

\[
| \mathrm{G} | ^ {- 1} \sum_ {g \in \mathrm{G}} \operatorname{Tr} (u \pi_ {\lambda} (g)) \operatorname{Tr} (v \pi_ {\lambda} (g ^ {- 1})) = d _ {\lambda} ^ {- 1} \operatorname{Tr} (u v).\tag{20}
\]

把这一关系式应用于 u 与 v 的秩 \(\leqslant 1\) 的情形；我们得到

\[
| \mathrm{G} | ^ {- 1} \sum_ {g \in \mathrm{G}} \left\langle x ^ {*}, \pi_ {\lambda} (g) x \right\rangle \left\langle y ^ {*}, \pi_ {\lambda} \left(g ^ {- 1}\right) y \right\rangle = d _ {\lambda} ^ {- 1} \left\langle x ^ {*}, y \right\rangle \left\langle y ^ {*}, x \right\rangle\tag{21}
\]

其中 \(x, y\) 属于 \(\mathrm{V}_{\lambda}\)，\(x^{*}, y^{*}\) 属于对偶 \(\mathrm{V}_{\lambda}^{*}\)（\(\mathrm{V}_{\lambda}\) 的对偶）。

对每个 \(\lambda \in \widehat{\mathbf{G}}\)，设 \((e_{\lambda,j})_{1\leqslant j\leqslant d_{\lambda}}\) 是 \(\mathrm{V}_{\lambda}\) 的一组基；用 \((\pi_{ij}^{\lambda}(g))\) 表示自同态 \(\pi_{\lambda}(g)\)（\(\mathrm{V}_{\lambda}\) 的自同态）关于这组基的矩阵。若用 \((e_{\lambda,i}^{*})_{1\leqslant i\leqslant d_{\lambda}}\) 表示 \(\mathrm{V}_{\lambda}^{*}\) 的与 \((e_{\lambda,j})\) 对偶的基，则有 \(\pi_{ij}^{\lambda}(g) = \langle e_{\lambda,i}^{*}, \pi_{\lambda}(g)e_{\lambda,j}\rangle\)，因此

\[
| \mathrm{G} | ^ {- 1} \sum_ {g \in \mathrm{G}} \pi_ {i j} ^ {\lambda} (g) \pi_ {k \ell} ^ {\lambda} (g ^ {- 1}) = d _ {\lambda} ^ {- 1} \delta_ {i \ell} \delta_ {j k}.\tag{22}
\]

现设 \(\lambda\) 与 \(\mu\) 是 \(\widehat{\mathbf{G}}\) 的两个不同元素，且 \(u \in \mathrm{End}_{\mathrm{K}}(\mathrm{V}_{\lambda})\)，\(v \in \mathrm{End}_{\mathrm{K}}(\mathrm{V}_{\mu})\)。同样由关系式 (15)，我们有

\[
\sum_ {g \in \mathrm{G}} \operatorname{Tr} (u \pi_ {\lambda} (g)) \operatorname{Tr} (v \pi_ {\mu} (g ^ {- 1})) = 0.\tag{23}
\]

如上所述，我们导出

\[
\sum_ {g \in \mathrm{G}} \langle x ^ {*}, \pi_ {\lambda} (g) x \rangle \langle y ^ {*}, \pi_ {\mu} (g ^ {- 1}) y \rangle = 0\tag{24}
\]

其中 \(x\in \mathrm{V}_{\lambda},x^{*}\in \mathrm{V}_{\lambda}^{*},y\in \mathrm{V}_{\mu}\)，以及 \(y^{*}\in \mathrm{V}_{\mu}^{*}\)。我们还有

\[
\sum_ {g \in \mathrm{G}} \pi_ {i j} ^ {\lambda} (g) \pi_ {k \ell} ^ {\mu} (g ^ {- 1}) = 0\tag{25}
\]

对 \(i,j\)（在 \([1,d_{\lambda}]\) 中）与 \(k,\ell\)（在 \([1,d_{\mu}]\) 中）成立。

关系式 (20) 至 (25) 称为舒尔正交关系。

注. —— 我们把代数 \(\operatorname{End}_{\mathrm{K}}(\mathrm{V}_{\lambda})\) 与矩阵代数 \(\mathbf{M}_{d_{\lambda}}(\mathrm{K})\) 等同，借助基 \((e_{\lambda,j})\)（\(V_{\lambda}\) 的基）。映射 \(\overline{F}^{-1}\) 是从代数 \(\prod_{\lambda}\mathbf{M}_{d_{\lambda}}(\mathrm{K})\) 到代数 K{[}G{]} 的一个同构。对 \(\mu\in\widehat{G}\)，用 \(E_{ij}^{\mu}\) 表示 \(\prod_{\lambda}\mathbf{M}_{d_{\lambda}}(\mathrm{K})\) 中这样的元素：其指标 \(\mu\) 的分量是矩阵单位 \(E_{ij}\)（\(\mathbf{M}_{d_{\mu}}(\mathrm{K})\) 的矩阵单位）（II, §10, No.~3, 第 341 页），其余分量均为零；令 \(u_{ij}^{\mu}=\overline{\mathcal{F}}^{-1}(\mathrm{E}_{ij}^{\mu})\)。元素族 \(u_{ij}^{\lambda}\)（\(\lambda\in\widehat{G}\)，\(1\leqslant i\leqslant d_{\lambda}\)，\(1\leqslant j\leqslant d_{\lambda}\)）是代数 K{[}G{]} 的一组基；其乘法表为

\[
u _ {i j} ^ {\lambda} u _ {k \ell} ^ {\mu} = \delta_ {\lambda \mu} \delta_ {j k} u _ {i \ell} ^ {\lambda}.\tag{26}
\]

此外，由公式 (15)，我们有

\[
u _ {i j} ^ {\lambda} = | \mathrm{G} | ^ {- 1} d _ {\lambda} \sum_ {g \in \mathrm{G}} \pi_ {j i} ^ {\lambda} (g ^ {- 1}) g.\tag{27}
\]

\subsection{特征标的正交关系}

我们保留第 5、6 小节的记号。回顾 \(Z(K[G])\) 由中心函数组成。

我们用公式定义一个从 \(K[G] \times K[G]\) 到 K 的对称双线性映射：

\[
\langle f, f ^ {\prime} \rangle_ {\mathrm{G}} = | \mathrm{G} | ^ {- 1} \sum_ {g \in \mathrm{G}} f _ {g} f _ {g ^ {- 1}} ^ {\prime}\tag{28}
\]

对每个 \(f = \sum f_{g}g\) 与 \(f' = \sum f_{g}'g\)（都属于 K{[}G{]}）成立。我们有 \(\langle f, f' \rangle_{\mathrm{G}} = |\mathrm{G}|^{-2}\tau (ff')\)。

命题 4（特征标的正交关系）

对 \(\lambda\) 与 \(\mu\)（在 \(\hat{\mathrm{G}}\) 中），我们有 \(\langle \chi_{\lambda},\chi_{\mu}\rangle_{\mathrm{G}} = \delta_{\lambda \mu}\)。

这是关系式 (20) 与 (23) 当自同态 u 与 v 取为恒同时的特殊情形。

推论. —— 设 \(\pi\) 与 \(\pi'\) 是 G 的有限维线性表示。在域 K 中，我们有

\[
\langle \chi_ {\pi}, \chi_ {\pi^ {\prime}} \rangle_ {\mathrm{G}} = (\dim_ {\mathrm{K}} \operatorname{Hom} _ {\mathrm{G}} (\pi , \pi^ {\prime})) \cdot 1.\tag{29}
\]

我们首先设 \(\pi\) 与 \(\pi'\) 是单表示。向量空间 \(\mathrm{Hom}_{\mathrm{G}}(\pi, \pi')\) 的维数依 \(\pi\) 与 \(\pi'\) 是否同构而为 1 或 0（舒尔引理，VIII, 第 47 页，命题 2）。在这一情形，公式 (29) 由命题 4 得出。

在一般情形，表示 \(\pi\)（相应地，\(\pi'\)）是单表示 \(\pi_{1},\ldots,\pi_{m}\)（相应地，\(\pi_{1}',\ldots,\pi_{n}'\)）的直和。空间 \(\operatorname{Hom}_{\mathrm{G}}(\pi,\pi')\) 同构于诸空间 \(\operatorname{Hom}_{\mathrm{G}}(\pi_{i},\pi_{j}')\)（\(1\leqslant i\leqslant m\)，\(1\leqslant j\leqslant n\)）的直和，并且我们有

\[
\chi_ {\pi} = \chi_ {\pi_ {1}} + \dots + \chi_ {\pi_ {m}}, \quad \chi_ {\pi^ {\prime}} = \chi_ {\pi_ {1} ^ {\prime}} + \dots + \chi_ {\pi_ {n} ^ {\prime}}.
\]

由线性性，公式 \((29)\) 的证明可化为单表示的情形。

\subsection{有限群上的中心函数}

命题 5. —— 族 \((\chi_{\lambda})_{\lambda\in\widehat{\mathbf{G}}}\) 是中心函数向量空间的一组基。G 的单线性表示的类的个数等于 G 的共轭类的个数。

对 \(a = \sum_{g \in G} a_g g \in K[G]\)，记 \(a^\vee = \sum_{g \in G} a_{g^{-1}} g\)；映射 \(a \mapsto a^\vee\) 是代数 K{[}G{]} 的一个对合反自同构。由公式 (17)，我们有 \(e_\lambda = |G|^{-1} d_\lambda \chi_\lambda^\vee\)（对每个 \(\lambda \in \widehat{G}\)）。于是该命题由族 \((e_\lambda)_{\lambda \in \widehat{G}}\) 是 K{[}G{]} 的中心的一组基这一事实得出。

用 C 表示 G 的共轭类所成的集合。设 C 是 C 的一个元素；像 \(C^{-1}\)（C 在映射 \(g \mapsto g^{-1}\) 下的像）是一个共轭类。C 中元素的交换子是 G 的彼此共轭的子群；其基数 \(d(\mathrm{C})\) 满足

\[
\operatorname{Card} (\mathrm{G}) = \operatorname{Card} (\mathrm{C}) d (\mathrm{C}).\tag{30}
\]

特别地，在域 K 中有 \(d(\mathrm{C}) \cdot 1 \neq 0\)。

设 f 是 G 上的一个中心函数。对任意共轭类 C，用 \(f(\mathrm{C})\) 表示诸 \(f(x)\)（\(x \in C\)）的公共值。用这一记号，特征标的正交关系（VIII, 第 410 页，命题 4）可写为

\[
\sum_ {\mathrm{C} \in \mathcal {C}} \chi_ {\lambda} (\mathrm{C} ^ {- 1}) \chi_ {\mu} (\mathrm{C}) d (\mathrm{C}) ^ {- 1} = \delta_ {\lambda \mu}\tag{31}
\]

对 \(\lambda\) 与 \(\mu\)（在 \(\widehat{G}\) 中）成立。

用 \(A\) 表示类型 \(\widehat{\mathbf{G}} \times \mathcal{C}\) 的、元素为 \(\chi_{\lambda}(\mathrm{C})\) 的矩阵，用 \(B\) 表示类型 \(\mathcal{C} \times \widehat{\mathbf{G}}\) 的、元素为 \(\chi_{\lambda}(\mathrm{C}^{-1}) d(\mathrm{C})^{-1}\) 的矩阵。集合 \(\widehat{\mathbf{G}}\) 与 \(\mathcal{C}\) 有相同的基数（命题 5）；关系式 (31) 表达了矩阵乘积 \(AB\) 是类型 \(\widehat{\mathbf{G}} \times \widehat{\mathbf{G}}\) 的单位矩阵这一事实。由 II, §10, No.~12, 第 360 页的命题 11，矩阵乘积 \(BA\) 是类型 \(\mathcal{C} \times \mathcal{C}\) 的单位矩阵；换句话说，我们有关系式

\[
\sum_ {\lambda \in \widehat {\mathrm{G}}} \chi_ {\lambda} (\mathrm{C} ^ {- 1}) \chi_ {\lambda} (\mathrm{C} ^ {\prime}) = d (\mathrm{C}) \delta_ {\mathrm{CC} ^ {\prime}}\tag{32}
\]

对 C 中的 C 与 C' 成立（有时称为"特征标的第二正交关系"）。

设 H 是 G 的一个子群。注意整数 \(\text{Card(H)}\) 整除 \(\text{Card(G)}\)，且 \(|G|\) 在 K 中不为零；因此在 K 中有 \(|H| \neq 0\)。

用 \(Res_{H}^{G}\) 表示从 \(Z(K[G])\) 到 \(Z(K[H])\) 的、把 G 上的中心函数映为其在 H 上的限制的线性映射。我们已经看到，若 \(\chi_{\pi}\) 是 G 的有限维表示 \(\pi\) 的特征标，则 \(\mathrm{Res}_{\mathrm{H}}^{\mathrm{G}}(\chi_{\pi})\) 是 H 的表示 \(\mathrm{Res}_{\mathrm{H}}^{\mathrm{G}}(\pi)\) 的特征标。

命题 6. —— 设 f 是 G 上的一个中心函数，u 是 H 上的一个中心函数。我们有

\[
\langle \mathrm{Ind} _ {\mathrm{H}} ^ {\mathrm{G}} (u), f \rangle_ {\mathrm{G}} = \langle u, \mathrm{Res} _ {\mathrm{H}} ^ {\mathrm{G}} (f) \rangle_ {\mathrm{H}}.\tag{33}
\]

G 的单表示的特征标构成 \(Z(\mathrm{K}[\mathrm{G}])\) 的一组基（VIII, 第 411 页，命题 5），对 H 亦然。因此只需在 f 是特征标 \(\chi_{\pi}\)（G 的单表示 \(\pi\) 的特征标）且 u 是特征标 \(\chi_{\sigma}\)（H 的单表示 \(\sigma\) 的特征标）的情形证明 (33)。在这一情形，\(\operatorname{Ind}_{\mathrm{H}}^{\mathrm{G}}(u)\) 是 G 的表示 \(\operatorname{Ind}_{\mathrm{H}}^{\mathrm{G}}(\sigma)\) 的特征标，并且由 VIII, 第 410 页的推论，我们有

\[
\langle \mathrm{Ind} _ {\mathrm{H}} ^ {\mathrm{G}} (u), f \rangle_ {\mathrm{G}} = (\dim_ {\mathrm{K}} \mathrm{Hom} _ {\mathrm{G}} (\mathrm{Ind} _ {\mathrm{H}} ^ {\mathrm{G}} (\sigma), \pi)) \cdot 1.
\]

我们同理可证关系式

\[
\langle u, \mathrm{Res} _ {\mathrm{H}} ^ {\mathrm{G}} (f) \rangle_ {\mathrm{H}} = (\dim_ {\mathrm{K}} \mathrm{Hom} _ {\mathrm{H}} (\sigma , \mathrm{Res} _ {\mathrm{H}} ^ {\mathrm{G}} (\pi))) \cdot 1
\]

，而等式 (33) 由弗罗贝尼乌斯互反律得出。

命题 7. —— 设 \(f\) 是从 \(G\) 到 \(K\) 的一个映射。下列断言等价：

\begin{enumerate}
\def\labelenumi{(\roman{enumi})}
\item
  存在元素 \(\lambda\)（\(\widehat{\mathbf{G}}\) 的元素）与元素 \(a\)（\(\mathrm{K}\) 的元素），使得 \(f = a\chi_{\lambda}\)。
\item
  对元素对 \((g,g^{\prime})\)（\(\mathbf{G}\) 的每一对元素），我们有
\end{enumerate}

\[
f (g) f (g ^ {\prime}) = | \mathrm{G} | ^ {- 1} f (1) \sum_ {h \in \mathrm{G}} f (h g h ^ {- 1} g ^ {\prime}).\tag{34}
\]

设 \(\lambda\) 是 \(\widehat{\mathbf{G}}\) 的一个元素；对每个自同态 \(u\)（\(\mathrm{V}_{\lambda}\) 的自同态），令

\[
u ^ {\natural} = | \mathrm{G} | ^ {- 1} \sum_ {h \in \mathrm{G}} \pi_ {\lambda} (h) u \pi_ {\lambda} (h ^ {- 1}).\tag{35}
\]

自同态 \(u^{\natural}\)（\(V_{\lambda}\) 的自同态）是 K{[}G{]}-线性的。由舒尔引理（VIII, 第 47 页，定理 1），\(u^{\natural}\) 是一个位似。由于 \(u\) 与 \(u^{\natural}\) 有相同的迹，因此我们有

\[
u ^ {\natural} = d _ {\lambda} ^ {- 1} \mathrm{Tr} (u) 1 _ {\mathrm{V} _ {\lambda}}.
\]

设 \(u\) 与 \(v\) 是 \(V_{\lambda}\) 的自同态；由此得出

\[
\operatorname{Tr} (u) \operatorname{Tr} (v) = d _ {\lambda} \operatorname{Tr} (u ^ {\natural} v) = d _ {\lambda} | G | ^ {- 1} \sum_ {h \in G} \operatorname{Tr} (\pi_ {\lambda} (h) u \pi_ {\lambda} (h ^ {- 1}) v).\tag{36}
\]

在公式 (36) 中取 \(u = \pi_{\lambda}(g)\) 与 \(v = \pi_{\lambda}(g')\)；即得 \(f = \chi_{\lambda}\) 时的关系式 (34)。

反之，设断言 (ii) 成立。若 \(f(1)=0\)，则 \(f(g)f(g')=0\) 对 G 的每一对元素 \((g,g')\) 成立，从而 f=0。因此可设 \(f(1)\neq0\)。在 (34) 中取 \(g'=1\)，我们得到关系式

\[
f (g) = | \mathrm{G} | ^ {- 1} \sum_ {h \in \mathrm{G}} f (h g h ^ {- 1})
\]

对每个 \(g \in G\) 成立，这蕴含 f 是一个中心函数。由命题 5，存在 K 的元素的一个族 \((a_{\lambda})\)，使得 \(f = \sum_{\lambda \in \widehat{G}} a_{\lambda} \chi_{\lambda}\)。把这个表达式代入公式 (34) 中的 f；注意到每个特征标 \(\chi_{\lambda}\) 也满足这一关系式，我们得到

\[
\sum_ {\lambda , \mu \in \widehat {G}} a _ {\lambda}   a _ {\mu}   \chi_ {\lambda} (g)   \chi_ {\mu} (g ^ {\prime}) = \sum_ {\lambda \in \widehat {G}} a _ {\lambda}   d _ {\lambda} ^ {- 1}   f (1)   \chi_ {\lambda} (g)   \chi_ {\lambda} (g ^ {\prime})\tag{37}
\]

对 \(g, g' \in \mathbf{G}\) 成立。这一关系式也可写为

\[
\sum_ {\lambda , \mu} (a _ {\lambda} a _ {\mu} - \delta_ {\lambda \mu} a _ {\lambda} d _ {\lambda} ^ {- 1} f (1)) \chi_ {\lambda} (g) \chi_ {\mu} (g ^ {\prime}) = 0\tag{38}
\]

对 \(g, g' \in \mathrm{G}\) 成立。而今，诸函数 \(\chi_{\lambda}\)（\(\lambda \in \widehat{\mathrm{G}}\)）线性无关（VIII, 第 411 页，命题 5）；由此得出

\[
a _ {\lambda} a _ {\mu} = \delta_ {\lambda \mu} a _ {\lambda} d _ {\lambda} ^ {- 1} f (1)
\]

对 \(\lambda, \mu \in \widehat{\mathbf{G}}\) 成立。特别地，\(a_{\lambda}a_{\mu} = 0\)（只要 \(\lambda \neq \mu\)）。因此，至多存在一个元素 \(\lambda\)（\(\widehat{\mathbf{G}}\) 的元素）使 \(a_{\lambda} \neq 0\)，于是 \(f = a_{\lambda}\chi_{\lambda}\)，从而得到 (i)。

\subsection{交换群的情形}

在本小节中，我们假设群 \(\mathrm{G}\) 是交换的。

由舒尔引理（VIII, 第 48 页，推论 1），G 的每个单表示的维数都是 1。设 \((\mathrm{M},\pi)\) 是这样一个表示，\(\chi\) 是它的特征标；对每个 \(g\in G\) 与 \(x\in M\)，有 \(\pi(g)(x)=\chi(g)x\)。于是，特征标 \(\chi\) 是从 G 到 K 的乘法群 \(K^{*}\) 的一个同态。反之，从 G 到 \(K^{*}\) 的每个同态都是 G 的一个次数为 1 的表示的特征标。因此，集合 \(\widehat{G}\)（单 K{[}G{]}-模的类所成的集合）可与集合 \(\operatorname{Hom}(G,K^{*})\)（从 G 到 \(K^{*}\) 的同态所成的集合）等同。由此我们在 \(\widehat{G}\) 上导出一个交换群结构；\(\widehat{G}\) 中的乘积对应于表示的张量积。由 VIII, 第 411 页的命题 5，群 G 与 \(\widehat{G}\) 有相同的基数。G 上的每个函数都是中心函数，且 \(\widehat{G}\) 是从 G 到 K 的映射所成的向量空间的一组基（同前引文）。由特征标的正交关系，这样的映射 f 关于基 \(\widehat{G}\) 有如下唯一的分解：

\[
f = \sum_ {\chi \in \widehat {\mathrm{G}}} \langle \chi , f \rangle_ {\mathrm{G}} \chi .\tag{39}
\]

对映射 \(f\) 与 \(f'\)（从 \(G\) 到 \(K\) 的映射），我们有关系式

\[
\langle f, f ^ {\prime} \rangle_ {\mathrm{G}} = \sum_ {\chi \in \widehat {\mathrm{G}}} \langle \chi , f \rangle_ {\mathrm{G}} \langle \chi , f ^ {\prime} \rangle_ {\mathrm{G}}.\tag{40}
\]

设 \((\mathrm{V},\pi)\) 是 G 的一个线性表示。对任意 \(\chi\in\widehat{\mathrm{G}}\)，用 \(V^{\chi}\) 表示 V 中由满足 \(\pi(g)(v)=\chi(g)v\)（对每个 \(g\in\mathrm{G}\)）的向量 v 组成的子空间。空间 \(V^{\chi}\) 是 \(\chi\) 型同型分量（K{[}G{]}-模 V 的同型分量）。空间 V 是族 \((V^{\chi})_{\chi\in\widehat{\mathrm{G}}}\) 的直和，而与这一分解相伴的 V 的投影算子 \(p_{\chi}\)（以 \(V^{\chi}\) 为像的投影算子）由

\[
p _ {\chi} = | \mathrm{G} | ^ {- 1} \sum_ {g \in \mathrm{G}} \chi (g ^ {- 1}) \pi (g)\tag{41}
\]

给出（依关系式 (19)）。

注. —— 设 n 是群 G 的基数，\(\mu_{n}(\mathrm{K})\) 是 K 中 n 次单位根所成的群。对每个 \(g \in G\)，有 \(g^{n} = 1\)；于是，\(\widehat{G}\) 可与群 \(\operatorname{Hom}(\mathbf{G}, \mu_{n}(\mathbf{K}))\) 等同。群 \(\mu_{n}(\mathbf{K})\) 是 n 阶循环群（V, §11, No.~2, 第 78 页，定理 1）。因此群 \(\widehat{G}\) 同构于群 \(\mathrm{D}(\mathrm{G}) = \operatorname{Hom}(\mathrm{G}, \mathbf{Q}/\mathbf{Z})\)。由 VII, §4, No.~9, 第 26 页的

命题 10，群 \(\widehat{\mathbf{G}}\) 同构于群 \(\mathbf{G}\)，且把元素 \(g\)（\(\mathbf{G}\) 的元素）映为同态 \(\chi \mapsto \chi(g)\)（从 \(\widehat{\mathbf{G}}\) 到 \(\mathrm{K}^*\) 的同态）的映射是从 \(\mathbf{G}\) 到 \(\widehat{\widehat{\mathbf{G}}}\) 的一个同构。
